% Options for packages loaded elsewhere
\PassOptionsToPackage{unicode}{hyperref}
\PassOptionsToPackage{hyphens}{url}
%
\documentclass[
]{article}
\usepackage{amsmath,amssymb}
\usepackage{iftex}
\ifPDFTeX
  \usepackage[T1]{fontenc}
  \usepackage[utf8]{inputenc}
  \usepackage{textcomp} % provide euro and other symbols
\else % if luatex or xetex
  \usepackage{unicode-math} % this also loads fontspec
  \defaultfontfeatures{Scale=MatchLowercase}
  \defaultfontfeatures[\rmfamily]{Ligatures=TeX,Scale=1}
\fi
\usepackage{lmodern}
\ifPDFTeX\else
  % xetex/luatex font selection
\fi
% Use upquote if available, for straight quotes in verbatim environments
\IfFileExists{upquote.sty}{\usepackage{upquote}}{}
\IfFileExists{microtype.sty}{% use microtype if available
  \usepackage[]{microtype}
  \UseMicrotypeSet[protrusion]{basicmath} % disable protrusion for tt fonts
}{}
\makeatletter
\@ifundefined{KOMAClassName}{% if non-KOMA class
  \IfFileExists{parskip.sty}{%
    \usepackage{parskip}
  }{% else
    \setlength{\parindent}{0pt}
    \setlength{\parskip}{6pt plus 2pt minus 1pt}}
}{% if KOMA class
  \KOMAoptions{parskip=half}}
\makeatother
\usepackage{xcolor}
\usepackage{color}
\usepackage{fancyvrb}
\newcommand{\VerbBar}{|}
\newcommand{\VERB}{\Verb[commandchars=\\\{\}]}
\DefineVerbatimEnvironment{Highlighting}{Verbatim}{commandchars=\\\{\}}
% Add ',fontsize=\small' for more characters per line
\newenvironment{Shaded}{}{}
\newcommand{\AlertTok}[1]{\textcolor[rgb]{1.00,0.00,0.00}{\textbf{#1}}}
\newcommand{\AnnotationTok}[1]{\textcolor[rgb]{0.38,0.63,0.69}{\textbf{\textit{#1}}}}
\newcommand{\AttributeTok}[1]{\textcolor[rgb]{0.49,0.56,0.16}{#1}}
\newcommand{\BaseNTok}[1]{\textcolor[rgb]{0.25,0.63,0.44}{#1}}
\newcommand{\BuiltInTok}[1]{\textcolor[rgb]{0.00,0.50,0.00}{#1}}
\newcommand{\CharTok}[1]{\textcolor[rgb]{0.25,0.44,0.63}{#1}}
\newcommand{\CommentTok}[1]{\textcolor[rgb]{0.38,0.63,0.69}{\textit{#1}}}
\newcommand{\CommentVarTok}[1]{\textcolor[rgb]{0.38,0.63,0.69}{\textbf{\textit{#1}}}}
\newcommand{\ConstantTok}[1]{\textcolor[rgb]{0.53,0.00,0.00}{#1}}
\newcommand{\ControlFlowTok}[1]{\textcolor[rgb]{0.00,0.44,0.13}{\textbf{#1}}}
\newcommand{\DataTypeTok}[1]{\textcolor[rgb]{0.56,0.13,0.00}{#1}}
\newcommand{\DecValTok}[1]{\textcolor[rgb]{0.25,0.63,0.44}{#1}}
\newcommand{\DocumentationTok}[1]{\textcolor[rgb]{0.73,0.13,0.13}{\textit{#1}}}
\newcommand{\ErrorTok}[1]{\textcolor[rgb]{1.00,0.00,0.00}{\textbf{#1}}}
\newcommand{\ExtensionTok}[1]{#1}
\newcommand{\FloatTok}[1]{\textcolor[rgb]{0.25,0.63,0.44}{#1}}
\newcommand{\FunctionTok}[1]{\textcolor[rgb]{0.02,0.16,0.49}{#1}}
\newcommand{\ImportTok}[1]{\textcolor[rgb]{0.00,0.50,0.00}{\textbf{#1}}}
\newcommand{\InformationTok}[1]{\textcolor[rgb]{0.38,0.63,0.69}{\textbf{\textit{#1}}}}
\newcommand{\KeywordTok}[1]{\textcolor[rgb]{0.00,0.44,0.13}{\textbf{#1}}}
\newcommand{\NormalTok}[1]{#1}
\newcommand{\OperatorTok}[1]{\textcolor[rgb]{0.40,0.40,0.40}{#1}}
\newcommand{\OtherTok}[1]{\textcolor[rgb]{0.00,0.44,0.13}{#1}}
\newcommand{\PreprocessorTok}[1]{\textcolor[rgb]{0.74,0.48,0.00}{#1}}
\newcommand{\RegionMarkerTok}[1]{#1}
\newcommand{\SpecialCharTok}[1]{\textcolor[rgb]{0.25,0.44,0.63}{#1}}
\newcommand{\SpecialStringTok}[1]{\textcolor[rgb]{0.73,0.40,0.53}{#1}}
\newcommand{\StringTok}[1]{\textcolor[rgb]{0.25,0.44,0.63}{#1}}
\newcommand{\VariableTok}[1]{\textcolor[rgb]{0.10,0.09,0.49}{#1}}
\newcommand{\VerbatimStringTok}[1]{\textcolor[rgb]{0.25,0.44,0.63}{#1}}
\newcommand{\WarningTok}[1]{\textcolor[rgb]{0.38,0.63,0.69}{\textbf{\textit{#1}}}}
\usepackage{graphicx}
\makeatletter
\def\maxwidth{\ifdim\Gin@nat@width>\linewidth\linewidth\else\Gin@nat@width\fi}
\def\maxheight{\ifdim\Gin@nat@height>\textheight\textheight\else\Gin@nat@height\fi}
\makeatother
% Scale images if necessary, so that they will not overflow the page
% margins by default, and it is still possible to overwrite the defaults
% using explicit options in \includegraphics[width, height, ...]{}
\setkeys{Gin}{width=\maxwidth,height=\maxheight,keepaspectratio}
% Set default figure placement to htbp
\makeatletter
\def\fps@figure{htbp}
\makeatother
\setlength{\emergencystretch}{3em} % prevent overfull lines
\providecommand{\tightlist}{%
  \setlength{\itemsep}{0pt}\setlength{\parskip}{0pt}}
\setcounter{secnumdepth}{-\maxdimen} % remove section numbering
\ifLuaTeX
\usepackage[bidi=basic]{babel}
\else
\usepackage[bidi=default]{babel}
\fi
\babelprovide[main,import]{russian}
% get rid of language-specific shorthands (see #6817):
\let\LanguageShortHands\languageshorthands
\def\languageshorthands#1{}
\ifLuaTeX
  \usepackage{selnolig}  % disable illegal ligatures
\fi
\usepackage{bookmark}
\IfFileExists{xurl.sty}{\usepackage{xurl}}{} % add URL line breaks if available
\urlstyle{same}
\hypersetup{
  pdftitle={AB-облако как фазовый резонатор},
  pdfauthor={Исаев Исхак Хамзатович (ORCID: 0009-0003-7299-0701)},
  pdflang={ru},
  pdfkeywords={гипотеза Гильберта--Поля, дзета-функция Римана, теория
случайных матриц, GUE, эффект Ааронова--Бома, гамильтониан
Хофштадтера, AB-облако, топологические вихри, дираковский конус, TKNN},
  hidelinks,
  pdfcreator={LaTeX via pandoc}}

\title{AB-облако как фазовый резонатор}
\usepackage{etoolbox}
\makeatletter
\providecommand{\subtitle}[1]{% add subtitle to \maketitle
  \apptocmd{\@title}{\par {\large #1 \par}}{}{}
}
\makeatother
\subtitle{Нули дзета-функции Римана, GUE-универсальность и
топологическое вихревое вещество: полное численное исследование с
независимой верификацией}
\author{Исаев Исхак Хамзатович (ORCID: 0009-0003-7299-0701)}
\date{Версия 22 · Август 2026}

\begin{document}
\maketitle

\section{Аннотация}\label{ux430ux43dux43dux43eux442ux430ux446ux438ux44f}

Настоящая монография написана заново, «с чистого листа», и подводит итог
систематического численного исследования объекта, который мы называем
\textbf{AB-облаком}: решёточной квантовой системы (гамильтониан
Хофштадтера с фазами Ааронова--Бома), в которой топологические вихревые
возбуждения создают неоднородное фазовое поле, а распределение
собственных значений сравнивается со статистикой нулей дзета-функции
Римана. Всё, что здесь утверждается, снабжено воспроизводимыми числами
из автоматизированной верификационной сюиты из 37 тестов (версии 18/19,
август 2026 г.); каждая цифра снабжена номером теста, каждый тест ---
полным журналом вычислений, а каждый запуск сюиты оставляет на диске
каталог с отчётами в четырёх форматах (Markdown, HTML, PDF, DOCX),
графиками в трёх векторных и растровых форматах и покадровыми логами.

Главные подтверждённые результаты таковы. Во-первых, в гладкой
калибровке (модель фаз \texttt{:monumental}) среднее отношение соседних
зазоров спектра AB-облака на торе достигает
\(\langle r\rangle = 0{,}585\pm0{,}026\) при целевом значении GUE
\(0{,}5992\) --- отклонение \(-2{,}4\%\) (тест 33), а сканирование по
размерам решёток \(L=10,20,30,50\) даёт монотонную сходимость
\(0{,}548\to0{,}592\to0{,}605\to0{,}595\) (тест 34). Во-вторых, прямое
сравнение развёрнутых спейсингов AB-облака и первых 5000 нулей
\(\zeta(s)\) обнаруживает, что парная корреляционная функция облака
ближе к GUE, чем к Пуассону, по всему интервалу \(s\in[0;3]\), хотя
двухвыборочный критерий Колмогорова--Смирнова на таком малом ансамбле
отвергает гипотезу совпадения распределений (\(p=6{,}7\cdot10^{-4}\);
тест 35) --- мы честно фиксируем это расхождение и объясняем его
конечномерностью ансамбля. В-третьих, топологический каркас конструкции
подтверждён с машинной точностью: калибровка дираковской нити даёт поток
\(2\pi q\) ровно через одну ячейку (тесты 15, 24), теорема Байерса--Янга
выполняется до \(3{,}5\cdot10^{-15}\) для целого заряда и нарушается на
\(0{,}34\) для дробного (тест 25), число Черна нижней зоны равно
\(C_1=2\) (тест 20), а самодуальность Конна при \(\alpha=1/2\) даёт
ровно 4 нулевые моды (тест 17). Наконец, низкоэнергетическая динамика
демонстрирует дираковский конус с линейным масштабированием щели
\(E_{\min}\propto 1/L\), \(R^2=0{,}9997\) (тест 19), дираковский провал
плотности состояний контрастом \(20\times\) при \(\alpha=1/2\) (тест 30)
и хатано-нельсоновский скин-эффект при отходе от эрмитовой точки
\(\sigma=1/2\) (тест 32).

Совокупность результатов образует связный нарратив: AB-облако --- это
фазовый резонатор, в котором универсальная GUE-статистика возникает
из-за нарушения симметрии обращения времени комплексными фазами прыжков,
вихри ведут себя как релятивистские частицы, а критическая прямая
\(\sigma=1/2\) выделяется как точка самодуальности, в которой статистика
облака ближе всего к статистике нулей дзета-функции. Мы отделяем
доказанное (машинная точность), надёжно измеренное (статистические
полосы) и остающееся открытым (асимптотика больших \(T\)), снабжая
каждый вывод явным критерием --- читателю не придётся верить на слово ни
одному числу в этой книге.

\textbf{Ключевые слова:} гипотеза Гильберта--Поля; дзета-функция Римана;
теория случайных матриц; GUE; эффект Ааронова--Бома; гамильтониан
Хофштадтера; AB-облако; топологические вихри; дираковский конус; TKNN;
скин-эффект Хатано--Нельсона.

\section{1. Введение и постановка
задачи}\label{ux432ux432ux435ux434ux435ux43dux438ux435-ux438-ux43fux43eux441ux442ux430ux43dux43eux432ux43aux430-ux437ux430ux434ux430ux447ux438}

\subsection{1.1 Гипотеза
Гильберта--Поля}\label{ux433ux438ux43fux43eux442ux435ux437ux430-ux433ux438ux43bux44cux431ux435ux440ux442ux430ux43fux43eux43bux44f}

В 1914 году Риман высказал гипотезу, что все нетривиальные нули
дзета-функции \(\zeta(s)\) лежат на критической прямой
\(\mathrm{Re}\,s = 1/2\). Спустя почти полвека, в 1958 году, Дэвид
Гильберт (по воспоминаниям Пойи) предложил стратегию доказательства,
которая с тех пор называется \textbf{гипотезой Гильберта--Поля}: если
найти самосопряжённый оператор \(\hat H\), спектр которого совпадает с
ординатами нулей \(\{\gamma_n\}\), то гипотеза Римана стала бы теоремой
--- эрмитовость оператора немедленно влечёт вещественность всех
собственных значений, то есть \(\mathrm{Re}\,\gamma_n = 1/2\) для всех
\(n\). Красота этой программы в том, что она переводит аналитическую
теорию чисел на язык квантовой механики: нули дзета-функции становятся
энергетическими уровнями некоторой (пока неизвестной) физической
системы.

Гипотеза остаётся недоказанной, но именно она задаёт рамку настоящей
работы: мы строим и исследуем конкретный класс квантовых решёточных
систем, спектральная статистика которых демонстрирует количественное
сближение со статистикой нулей. Важно подчеркнуть методологическую
позицию, отличающую эту монографию от многих предшествующих текстов на
данную тему: мы \textbf{не утверждаем доказательства гипотезы Римана}.
Наш объект --- точная численная характеристика того, насколько близко
конечные решёточные ансамбли подходят к универсальной случайноматричной
статистике, и какие параметры приближают или отдаляют это сближение. Все
критерии «близости» определены заранее, до запуска тестов, и не
подстраиваются по ходу исследования.

Исторический контекст коротко таков. В 1972 году Монтгомери
сформулировал гипотезу о парной корреляции нулей \(\zeta\); Дисон,
случайно встретившись с Монтгомери в чаепательной комнате Принстонского
института перспективных исследований, указал, что формула Монтгомери
совпадает с парной корреляцией собственных значений ансамбля GUE
(гауссов унитарный ансамбль) теории случайных матриц. В 1980-е годы
Одлыжко на огромном вычислительном материале подтвердил это совпадение
численно, а Берри в 1985 году вычислил конечномерные поправки: реальная
(ограниченная по \(T\)) выборка нулей должна отклоняться от чистого GUE
предсказуемым образом --- в частности, значение парной корреляции в нуле
\(R_2(0)\) получает поправку, экспоненциально или степенно малую по
\(T\). Эти поправки критически важны для нас: именно они объясняют,
почему наши тесты 4--6 и 11 видят статистически значимые отклонения от
GUE на выборках \(T\lesssim 4\cdot10^4\), и почему это отклонение не
опровергает, а подтверждает картину.

\subsection{1.2 Эффект Ааронова--Бома и
AB-облако}\label{ux44dux444ux444ux435ux43aux442-ux430ux430ux440ux43eux43dux43eux432ux430ux431ux43eux43cux430-ux438-ab-ux43eux431ux43bux430ux43aux43e}

\textbf{Эффект Ааронова--Бома (АБ)} состоит в том, что квантовая частица
реагирует на магнитный поток даже там, где напряжённость поля равна
нулю: волновая функция, огибающая соленоид с потоком \(\Phi\),
приобретает фазу \(\varphi_{\rm AB} = 2\pi\Phi/\Phi_0\), где
\(\Phi_0 = h/e\) --- квант потока. На решётке этот эффект кодируется
заменой амплитуд прыжков: \(t_{ij}\to t_{ij}e^{i\varphi_{ij}}\), где
фазы определяются векторным потенциалом. \textbf{Гамильтониан
Хофштадтера} --- это частный случай такой конструкции: электроны на
квадратной решётке в однородном магнитном поле с потоком
\(\alpha = \Phi/\Phi_0\) на ячейку; его спектр как функция \(\alpha\)
образует знаменитую «бабочку Хофштадтера» (она воспроизведена на рис. 20
нашего набора графиков).

\textbf{AB-облаком} мы называем обобщение этой модели, в котором на
решётку помещены точечные \textbf{топологические вихри} --- центры,
вокруг которых фаза прыжков закручена на \(2\pi q\) (\(q\) --- заряд
вихря, может быть дробным). Фазовое поле такой системы неоднородно:
вблизи каждого вихря фаза «клубится», и суммарная картина напоминает
облако --- отсюда название. Математически гамильтониан записывается как

\[H_{ij} = -\exp\bigl(i\varphi_{ij}\bigr),\qquad \varphi_{ij} = 2\pi\alpha\,j + \sum_{k=1}^{N_v} q_k\bigl[\arg(\mathbf r_i - \mathbf r_k) - \arg(\mathbf r_j - \mathbf r_k)\bigr] + \varphi^{\rm str}_{ij},\]

где \(j\) --- номер столбца решётки (калибровка Ландау), \(N_v\) ---
число вихрей, а \(\varphi^{\rm str}_{ij}\) --- явная \textbf{дираковская
нить}: добавка \(-2\pi q\), приписываемая тем связям, которые пересекают
полупрямую (нить), выходящую из вихря. Дираковская нить гарантирует, что
суммарный магнитный поток вокруг вихря равен ровно \(2\pi q\) с
точностью до машинного эпсилона --- это проверено прямыми вычислениями
(тесты 15, 24).

Кроме геометрической (фазовой) части, модель содержит
\textbf{диагональный беспорядок} \(W\): аддитивный шум на узлах,
моделирующий примеси реального кристалла. Параметры модели --- размер
решётки \(L\times L\), поток \(\alpha\), набор зарядов \(\{q_k\}\),
число и позиции вихрей \(N_v\), сила беспорядка \(W\) и тип граничных
условий (открытые или тор). Ключевой физический переключатель ---
\textbf{модель фаз вихрей}: (а) калибровка дираковской нити, в которой
целозарядные вихри спектрально невидимы (это точная теорема
Байерса--Янга, тест 25), и (б) гладкая калибровка \texttt{:monumental}
(atan-модель), в которой фазы распределены плавно и комплексны уже при
любом \(q\) --- именно она делает GUE-статистику достижимой при
физически разумных параметрах. Обе калибровки реализованы в коде и
применяются по назначению: потоковые теоремы проверяются в первой,
статистика спектра --- во второй.

\subsection{1.3 Связь с теорией случайных
матриц}\label{ux441ux432ux44fux437ux44c-ux441-ux442ux435ux43eux440ux438ux435ux439-ux441ux43bux443ux447ux430ux439ux43dux44bux445-ux43cux430ux442ux440ux438ux446}

\textbf{Теория случайных матриц (ТСМ)} изучает распределения собственных
значений больших матриц со случайными элементами. Для нас существенны
три ансамбля, различающиеся симметрией обращения времени: \textbf{GOE}
(вещественные симметричные матрицы, T-симметрия сохранена), \textbf{GUE}
(комплексные эрмитовы, T-симметрия нарушена) и \textbf{GSE}
(кватернионные самодвойственные), а также предельный
\textbf{пуассоновский} случай (некоррелированные уровни, интегрируемая
система). Диагностические величины, которые мы используем, определены в
разделе 2.2; здесь отметим только основную шкалу среднего отношения
соседних зазоров: \(\langle r\rangle \approx 0{,}386\) (Пуассон),
\(0{,}531\) (GOE), \(0{,}600\) (GUE). Тройная классификация Дайсона
предсказывает, что нарушение T-симметрии комплексными фазами обязано
переводить систему в класс GUE --- это «тройной путь Дайсона», и
AB-облако в гладкой калибровке является его пряной иллюстрацией: при
\(\alpha=1/2\) матрица комплексно-эрмитова,
\(\tau(H) = \|H-H^*\|/\|H\| = 0{,}225\neq 0\) (тест 26), и статистика
действительно тяготеет к GUE, тогда как в вещественной калибровке
дираковских нитей \(\tau = 0\) и спектр заперт в классе GOE (тесты 15,
16).

\subsection{1.4 Цели, задачи и организация
верификации}\label{ux446ux435ux43bux438-ux437ux430ux434ux430ux447ux438-ux438-ux43eux440ux433ux430ux43dux438ux437ux430ux446ux438ux44f-ux432ux435ux440ux438ux444ux438ux43aux430ux446ux438ux438}

Цель работы --- построить полную, независимо проверяемую картину
спектральной физики AB-облака. Задачи конкретизированы в виде 37
автоматизированных тестов, сгруппированных в шесть блоков: (1)
сходимость поправки \(b(N)\) и возражения критиков (тесты 1--14); (2)
построение и топологические теоремы (15, 17, 20, 24--27); (3)
GUE-классификация облака (16, 26, 29, 33, 34); (4) прямое сравнение с
нулями \(\zeta\) (28, 35, 36); (5) дираковская динамика и неэрмитовость
(19, 30--32, 37); (6) аналитические интерпретации (тесты 21--23 и раздел
8). Каждый тест возвращает булев вердикт по заранее зафиксированному
критерию; верификация двухпроходная (первый проход --- основная решётка,
второй --- увеличенная), а «быстрая проверка» прогоняет все 37 тестов на
решётках \(16\times16\to32\times32\) с усечённой до 5000 нулей
ζ-выборкой за считанные минуты.

Отдельно оговорим \textbf{две системы вердиктов}, чтобы читатель не
путал их. Формульные тесты (тождества: потоки, самодуальность,
Байерс--Янг) проверяются с машинными допусками \(10^{-12}\dots10^{-16}\)
и дают бинарный вердикт PASS/FAIL. Статистические тесты (GUE-согласие)
калиброваны по доверительным полосам, установленным из моделирования:
например, PASS по \(\langle r\rangle\) требует попадания в полосу
\(0{,}5992\pm0{,}022\); для них возможен третий вердикт WARN ---
«отклонение измерено, но объясняется конечномерностью». Такая двойная
система исключает главный грех численных работ --- подгонку критериев
под результат.

\section{2. Математический
аппарат}\label{ux43cux430ux442ux435ux43cux430ux442ux438ux447ux435ux441ux43aux438ux439-ux430ux43fux43fux430ux440ux430ux442}

\subsection{2.1 Дзета-функция Римана и
нули}\label{ux434ux437ux435ux442ux430-ux444ux443ux43dux43aux446ux438ux44f-ux440ux438ux43cux430ux43dux430-ux438-ux43dux443ux43bux438}

\textbf{Дзета-функция Римана} \(\zeta(s) = \sum_{n\ge1} n^{-s}\)
аналитически продолжается на всю комплексную плоскость и имеет
нетривиальные нули в критической полосе \(0<\mathrm{Re}\,s<1\). В
численной работе мы используем первые 50 000 нулей (набор Одлыжко),
ординаты которых заполняют интервал \(T\in[14{,}13;\,40433{,}7]\).
Плотность нулей на высоте \(T\) задаётся \textbf{формулой Римана--фон
Мангольдта}:

\[N(T) = \frac{T}{2\pi}\ln\frac{T}{2\pi} - \frac{T}{2\pi} + \frac{7}{8} + O(\ln T),\]

и растёт примерно как \((T/2\pi)\ln T\) --- нули «разрежаются» в
абсолютном расстоянии, но после деления на локальную плотность
становятся однородными. Эта процедура деления называется
\textbf{развёрткой (unfolding)}: нормированные спейсины
\(\delta_n = \gamma_{n+1}^{(u)} - \gamma_n^{(u)}\) имеют среднее 1 по
построению, и вся информация о корреляциях переносится в их
распределение. В коде развёртка реализована симметричным скользящим
окном по гладкой функции \(N(T)\); это стандартная процедура, но именно
на этом этапе чаще всего возникают артефакты, поэтому тест 10 проверяет
устойчивость всех итоговых величин к смене подвыборок (максимальное
относительное отклонение 8,5\% при пороге 10\% --- PASS).

Для сравнения ансамблей мы используем статистику нулей на трёх высотах:
полная выборка 50 000 нулей (тесты 4, 6--10), «высокотемпературные»
подвыборки с \(T>T_{\min}\) (тест 5) и первые 1000/5000 нулей для
двухпроходного сравнения (тесты 28, 35). Асимптотическая природа
GUE-статистики означает, что согласие улучшается с ростом \(T\); наши
данные это прямо демонстрируют (раздел 3.2).

\subsection{2.2 Ансамбли случайных матриц и диагностические
статистики}\label{ux430ux43dux441ux430ux43cux431ux43bux438-ux441ux43bux443ux447ux430ux439ux43dux44bux445-ux43cux430ux442ux440ux438ux446-ux438-ux434ux438ux430ux433ux43dux43eux441ux442ux438ux447ux435ux441ux43aux438ux435-ux441ux442ux430ux442ux438ux441ux442ux438ux43aux438}

Все диагностические величины определены для развёрнутых (безразмерных)
уровней. Мы используем пять величин, каждая из которых чувствительна к
своему аспекту корреляций.

\textbf{Отношение соседних зазоров} \(\langle r\rangle\). Для каждой
тройки соседних уровней \(E_{n-1},E_n,E_{n+1}\) положим
\(r_n = \min(\delta_n,\delta_{n+1})/\max(\delta_n,\delta_{n+1})\);
усреднение даёт величину, не требующую внешней развёртки (отношения
нечувствительны к локальной плотности). Опорные значения (большие
матрицы): Пуассон \(0{,}3863\), GOE \(0{,}5307\), GUE \(0{,}5996\). Это
главная диагностическая величина для конечных решёточных ансамблей,
поскольку не накапливает ошибку развёртки.

\textbf{Критерий Колмогорова--Смирнова (KS)}. Сравнивает эмпирическую
функцию распределения спейсингов с Wigner-предсказанием GUE: статистика
\(D\) = максимум модуля разности кривых, и при \(n\) спейсингах
критическое значение на уровне 5\% равно \(1{,}36/\sqrt{n}\). Для
\(n=49999\) критическое \(D = 0{,}0061\) --- то есть даже
микроскопические отклонения «отвергают» GUE; поэтому мы всегда приводим
\(D\) вместе с \(n\) и с критическим порогом (тест 4) и отдельно смотрим
тренд \(D(T_{\min})\) (тест 5).

\textbf{Числовая дисперсия \(\Sigma^2(L)\)} --- дисперсия числа уровней
в окне длины \(L\) (в единицах средней плотности). Пуассон даёт
\(\Sigma^2 = L\); GUE даёт
\(\Sigma^2 \approx (1/\pi^2)\ln(2\pi L) + \ldots\) --- логарифмический,
а не линейный рост. Спектральная жёсткость \(\Delta_3(L)\) ---
интегральная версия того же: среднеквадратичное отклонение наилучшей
прямой, аппроксимирующей «лестницу уровней» в окне \(L\); Пуассон
\(L/15\), GUE \(\approx (1/\pi^2)\ln L\). Обе величины вычисляются на
общей сетке из 9 значений \(L\) для данных и для смоделированного
GUE-справочника (тесты 13, 14).

\textbf{Парная корреляция \(R_2(s)\)} --- плотность вероятности найти
уровень на расстоянии \(s\) от данного (развитие гипотезы Монтгомери):
для GUE \(R_2(s) = 1 - \bigl(\sin \pi s/(\pi s)\bigr)^2\),
«корреляционная дыра» в нуле \(R_2(0)=0\) отражает выталкивание уровней;
для Пуассона \(R_2 \equiv 1\). Формфактор \(K(t)\) --- фурье-образ
\(R_2\), «рампа + плато» \(K(t)=\min(t,1)\) для GUE; это самая
интегральная из наших диагностик (тесты 35, 36).

\textbf{Критерий Андерсона--Дарлинга \(A^2\)} --- аналог KS, взвешенный
по хвостам распределения; p-значение мы оцениваем монте-карловски по
2000 симуляций (тест 11). Наконец, \textbf{f\_GUE} --- сводная мера
«доли GUE» в данных:
\(f_{\rm GUE} = (\Sigma^2_{\rm P} - \Sigma^2_{\rm data})/(\Sigma^2_{\rm P} - \Sigma^2_{\rm GUE})\)
по значению при \(L=2\); \(f_{\rm GUE}=1\) --- идеальный GUE, \(0\) ---
Пуассон (тест 29).

\subsection{2.3 Гамильтониан Хофштадтера с вихрями: точное
определение}\label{ux433ux430ux43cux438ux43bux44cux442ux43eux43dux438ux430ux43d-ux445ux43eux444ux448ux442ux430ux434ux442ux435ux440ux430-ux441-ux432ux438ux445ux440ux44fux43cux438-ux442ux43eux447ux43dux43eux435-ux43eux43fux440ux435ux434ux435ux43bux435ux43dux438ux435}

Соберём определение модели полностью. Решётка \(L_x\times L_y\);
волновая функция живёт на узлах; матричные элементы
\(H_{ij} = -\exp(i\varphi_{ij})\) для ближайших соседей (включая
периодические обходы, если выбран тор). Калибровка Ландау приписывает
столбцу \(j\) фазу \(2\pi\alpha j\). Вихрь с зарядом \(q_k\) в позиции
\(\mathbf r_k=(x_k,y_k)\) (координаты в ячейках) добавляет к разности
фаз связи \((i\to j)\) величину
\(q_k[\arg(\mathbf r_i-\mathbf r_k)-\arg(\mathbf r_j-\mathbf r_k)]\), а
дираковская нить --- дискретную добавку \(-2\pi q_k\) на связях,
пересекающих нить; в гладкой калибровке нить отсутствует, но atan-фаза
распределяется по вертикальным связям. Ключевое свойство: \textbf{сумма
фаз вокруг любой ячейки равна \(2\pi\alpha + 2\pi q_k\,\delta_k\)} (плюс
\(2\pi\cdot\)целое), где \(\delta_k\) --- индикатор того, что ячейка
содержит вихрь. Это проверено для всех 224 внутренних ячеек решётки
\(16\times16\): максимум \textbar поток\textbar{} на пустых ячейках
\(1{,}5\cdot10^{-14}\) рад (тест 24).

Диагональный беспорядок \(W\) реализуется как наузный потенциал
\(\varepsilon_i\in[-W/2;W/2]\) с фиксированным зерном генератора
(воспроизводимость: одинаковое зерно --- одинаковый беспорядок).
Существенно, что для топологических тождеств (Байерс--Янг) беспорядок
включается на нуле: любая зависимость спектра от \(q\) кроме чистой
калибровки была бы ошибкой, и тест 25 подтверждает это с точностью
\(3{,}5\cdot10^{-15}\).

Граничные условия --- отдельный физический выбор. Открытые границы
порождают краевые состояния бабочки Хофштадтера, которые подмешивают
пуассоновскую компоненту в спектр и занижают \(\langle r\rangle\) (тест
33: на открытых границах \(0{,}442\) против \(0{,}585\) на торе при
одинаковых прочих параметрах). Тор требует калибровочной совместности
\(q\,(2-N_x)\in\mathbb{Z}\) --- условие периодичности фазы при обходе;
для целых \(q\) оно выполняется автоматически, для дробных \(q=0{,}3\)
на торе \(16\times16\) оно нарушается, и система честно откатывается на
открытые границы (в коде это автоопределение с диагностическим
сообщением).

\subsection{2.4 Методология верификации и
воспроизводимость}\label{ux43cux435ux442ux43eux434ux43eux43bux43eux433ux438ux44f-ux432ux435ux440ux438ux444ux438ux43aux430ux446ux438ux438-ux438-ux432ux43eux441ux43fux440ux43eux438ux437ux432ux43eux434ux438ux43cux43eux441ux442ux44c}

Каждый запуск сюиты создаёт каталог
\texttt{results/run\_\textless{}дата-время\textgreater{}/} с
подкаталогом на каждый тест: внутри --- полный журнал вычислений (каждая
величина с меткой времени и дельтой относительно предыдущей), захват
консольного вывода, отчёты в форматах Markdown, HTML, PDF (векторный) и
DOCX, а также графики в PNG (600 dpi), SVG и PDF и покадровая
GIF-анимация первого графика. Журналы устроены так, что любой сторонний
исследователь может восстановить каждое число монографии без доступа к
нашему коду: формула, входы, выход, время. Команда воспроизведения
отдельного теста ---
\texttt{julia\ ab\_cloud\_v19.jl\ -\/-test\ \textless{}N\textgreater{}\ -\/-no-two-pass};
полной сюиты --- \texttt{-\/-test\ all}; быстрой проверки ---
\texttt{-\/-quick} (16×16→32×32, 5000 нулей, оба прохода). Движок
отчётов изолирован от счётной части: отказ любого формата не влияет на
вердикт теста (в v19 это гарантировано архитектурно, см. приложение A).

\section{\texorpdfstring{3. Блок 1: поправка \(b(N)\) и статистика нулей
\(\zeta\) (тесты
1--14)}{3. Блок 1: поправка b(N) и статистика нулей \textbackslash zeta (тесты 1--14)}}\label{ux431ux43bux43eux43a-1-ux43fux43eux43fux440ux430ux432ux43aux430-bn-ux438-ux441ux442ux430ux442ux438ux441ux442ux438ux43aux430-ux43dux443ux43bux435ux439-zeta-ux442ux435ux441ux442ux44b-114}

\subsection{\texorpdfstring{3.1 Что такое \(b(N)\) и почему её скорость
сходимости
важна}{3.1 Что такое b(N) и почему её скорость сходимости важна}}\label{ux447ux442ux43e-ux442ux430ux43aux43eux435-bn-ux438-ux43fux43eux447ux435ux43cux443-ux435ux451-ux441ux43aux43eux440ux43eux441ux442ux44c-ux441ux445ux43eux434ux438ux43cux43eux441ux442ux438-ux432ux430ux436ux43dux430}

Поправка \(b(N)\) возникает в аналитическом трактовке парной корреляции
нулей: это подгоночная величина, связывающая конечновыборочное значение
диагностик с их асимптотой. Ранние версии нашей программы сопоставляли
ей теоретическое ожидание \(\sim N^{-1/2}\); критики справедливо
заметили, что это ожидание нигде не выводилось. Мы принимаем это
возражение: начиная с версии 6 сюиты показатель степени \(b(N)\) ---
\textbf{чисто эмпирическая} величина без теоретического шаблона, и все
тесты блока 1 проверяют внутреннюю согласованность измерений, а не
соответствие несуществующей теореме. Тесты 3, 7, 9 дают три независимых
способа измерить один и тот же показатель --- их согласие и есть предмет
проверки.

\subsection{3.2 Три независимых измерения
показателя}\label{ux442ux440ux438-ux43dux435ux437ux430ux432ux438ux441ux438ux43cux44bux445-ux438ux437ux43cux435ux440ux435ux43dux438ux44f-ux43fux43eux43aux430ux437ux430ux442ux435ux43bux44f}

\begin{figure}
\centering
\includegraphics{../figures/fig01_bN_convergence.png}
\caption{\textbf{Рис. 1.} Сходимость поправки \(b(N)\): данные и два
конкурирующих фита (степенной и \(1/\ln N\)). Ожидание \(N^{-1/2}\), на
котором настаивали критики, данными исключается.}
\end{figure}

\textbf{Лог-лог регрессия} (тест 3) по сетке \(N\in[50;25600]\) даёт
\(b(N) \approx 7{,}03\,N^{-0{,}1685}\), \(R^2 = 0{,}9895\), стандартная
ошибка наклона \(0{,}0043\). \textbf{Расширенная регрессия до
\(N=32000\)} (тест 7) даёт наклон \(-0{,}1504\) с 95\%-м доверительным
интервалом \([-0{,}1594;-0{,}1414]\), \(R^2 = 0{,}9954\) --- интервалы
согласуются. \textbf{Бутстрап} (тест 9, 1000 перевыборок с возвратом)
даёт медиану наклона \(-0{,}1746\) и 95\%-й интервал
\([-0{,}1895;-0{,}1552]\). Важная деталь, которая обычно ускользает:
\textbf{альтернативная} параметризация
\(b(N) \approx -0{,}40 + 16{,}84/\ln N\) описывает данные лучше
степенного закона (\(R^2 = 0{,}9999\) против \(0{,}9954\) в тесте 7).
Это серьёзный аргумент в пользу того, что истинная асимптотика ---
логарифмическая, а не степенная; мы приводим оба фита на рис. 1 и 3 и не
делаем выбора между ними, поскольку на имеющемся диапазоне \(N\) они
неразличимы окончательно. Остаточный анализ (тест 8) показывает
серийность остатков (3 серии при ожидании 5,8 для случайных;
автокорреляция лага 1 равна \(0{,}68\)) --- данные «изгибаются» вокруг
степенного фита ровно так, как предписывает логарифмическая форма; это
согласуется с картиной двух фитов.

Практический вывод блока: на выборках \(N\le5\cdot10^4\) поправка
\(b(N)\) остаётся порядка единицы (от \(2{,}39\) при \(N=500\) до
\(1{,}28\) при \(N=32000\)) --- конечновыборочные эффекты в ζ-статистике
велики, и любой тест «GUE или нет», не учитывающий \(b(N)\), на этих
высотах просто некорректен.

\subsection{3.3 Возражение 2: KS, χ² и A² отвергают GUE --- и это
ожидаемо}\label{ux432ux43eux437ux440ux430ux436ux435ux43dux438ux435-2-ks-ux3c7uxb2-ux438-auxb2-ux43eux442ux432ux435ux440ux433ux430ux44eux442-gue-ux438-ux44dux442ux43e-ux43eux436ux438ux434ux430ux435ux43cux43e}

\begin{figure}
\centering
\includegraphics{../figures/fig02_spacing_hist.png}
\caption{\textbf{Рис. 2.} Гистограмма 49 999 спейсингов нулей \(\zeta\)
против сюрмизов Вигнера (GUE, GOE) и пуассоновской экспоненты: форма
типично «матричная», с отталкиванием уровней.}
\end{figure}

\begin{figure}
\centering
\includegraphics{../figures/fig05_ks_convergence.png}
\caption{\textbf{Рис. 3.} KS-статистика \(D\) против нижней границы
\(T_{\min}\): медленное убывание к критической линии --- численное
выражение асимптотичности GUE.}
\end{figure}

Тесты 4, 5, 6, 11 формулируют возражение критиков в его сильнейшей форме
и проверяют его честно. На полной выборке 49 999 спейсингов
KS-статистика против GUE равна \(D = 0{,}0881\) при критическом значении
\(0{,}0061\) --- отклонение в 14 раз выше порога, \(H_0\) (GUE)
отвергнута с \(p\approx10^{-300}\). Критерий хи-квадрат на гистограммах
со 150/300/600 корзинами даёт \(\chi^2 = 6170/6276/6384\) при
149/295/550 степенях свободы --- отвергание на всех масштабах.
Андерсон--Дарлинг на 5000 спейсингов против аналитической GUE-функции
распределения даёт \(A^2 = 117{,}7\) при монте-карловском 95-м
перцентиле \(16{,}6\) (\(p\approx0\) по 2000 симуляциям).

Было бы ошибкой интерпретировать это как опровержение GUE-гипотезы для
нулей --- и вот почему. GUE-статистика является
\textbf{асимптотической}: она достигается на высотах \(T\gg10^6\), где
накоплены миллионы нулей. Берри (1985) вычислил ведущие
конечновыборочные поправки, и их знак и порядок в точности соответствуют
наблюдаемым: на наших высотах \(T\le4\cdot10^4\) ожидаемо
\(D\sim0{,}08\). Прямое подтверждение дают тест 5 (высокотемпературный
срез: \(D = 0{,}0878\to0{,}0866\) при \(T_{\min}: 1000\to10000\) ---
медленное, но систематическое убывание) и, главное, тест 12:
\textbf{двухвыборочный} KS между реальными спейсингами и спейсингами
смоделированной GUE-матрицы \(2000\times2000\) даёт \(D = 0{,}0406\),
\(p = 0{,}139\) --- гипотеза «нулевые спейсинги распределены как
спейсинги GUE-матрицы того же размера» \textbf{не отвергается}. Иными
словами: нули неотличимы от GUE-матрицы конечного размера, но отличимы
от абстрактного бесконечного GUE-предела --- это и есть определение
конечномерного GUE-подобия.

\subsection{3.4 Жёсткость спектра: нули ближе к GUE, чем
GUE-моделирование}\label{ux436ux451ux441ux442ux43aux43eux441ux442ux44c-ux441ux43fux435ux43aux442ux440ux430-ux43dux443ux43bux438-ux431ux43bux438ux436ux435-ux43a-gue-ux447ux435ux43c-gue-ux43cux43eux434ux435ux43bux438ux440ux43eux432ux430ux43dux438ux435}

\begin{figure}
\centering
\includegraphics{../figures/fig03_decay_fits.png}
\caption{\textbf{Рис. 4.} Два фита скорости затухания \(b(N)\) в
лог-координатах: степенной и \(1/\ln N\); систематический изгиб остатков
указывает на логарифмическую асимптотику.}
\end{figure}

\begin{figure}
\centering
\includegraphics{../figures/fig04_bootstrap_slope.png}
\caption{\textbf{Рис. 5.} Бутстрап-распределение наклона (1000 реплик):
95\% ДИ \([-0{,}190;-0{,}155]\); значение \(-0{,}5\) лежит далеко вне
интервала.}
\end{figure}

\begin{figure}
\centering
\includegraphics{../figures/fig06_sigma2_L.png}
\caption{\textbf{Рис. 6.} Числовая дисперсия \(\Sigma^2(L)\): данные
(константа) против логарифмического GUE и линейного Пуассона --- 9/9
точек ближе к GUE.}
\end{figure}

\begin{figure}
\centering
\includegraphics{../figures/fig07_delta3_L.png}
\caption{\textbf{Рис. 7.} Спектральная жёсткость \(\Delta_3(L)\):
сверхжёсткость данных относительно даже смоделированного GUE на больших
\(L\).}
\end{figure}

Самые информативные результаты блока дают \(\Sigma^2(L)\) и
\(\Delta_3(L)\) (тесты 13, 14). На сетке из девяти значений
\(L\in[2;1400]\) данные нулей дают
\(\Sigma^2 \approx 0{,}33\text{–}0{,}41\) --- практически константу,
тогда как пуассоновское предсказание растёт линейно до 1400, а
смоделированный GUE-справочник --- логарифмически до
\(0{,}61\text{–}0{,}72\). Данные ближе к GUE, чем к Пуассону, на всех
9/9 точках; но они \textbf{жёстче} (меньше) даже смоделированного GUE на
больших \(L\) --- относительное превышение \(\Delta_3\) данных над
GUE-справочником падает от \(0{,}99\) при \(L=2\) до \(0{,}48\) при
\(L=1400\). Такая сверхжёсткость известна как особенность реальных нулей
\(\zeta\) (спектральная жёсткость арифметических систем выше матричной),
и её воспроизведение в наших данных --- самостоятельное подтверждение
чистоты развёртки: артефакты развёртки (сглаживание локальной плотности,
края окна) завышали бы \(\Delta_3\) и \(\Sigma^2\), а не занижали их.

\section{4. Блок 2: построение AB-облака и топологические теоремы (тесты
15, 17, 20,
24--27)}\label{ux431ux43bux43eux43a-2-ux43fux43eux441ux442ux440ux43eux435ux43dux438ux435-ab-ux43eux431ux43bux430ux43aux430-ux438-ux442ux43eux43fux43eux43bux43eux433ux438ux447ux435ux441ux43aux438ux435-ux442ux435ux43eux440ux435ux43cux44b-ux442ux435ux441ux442ux44b-15-17-20-2427}

\subsection{4.1 Корректность конструкции: эрмитовость и
потоки}\label{ux43aux43eux440ux440ux435ux43aux442ux43dux43eux441ux442ux44c-ux43aux43eux43dux441ux442ux440ux443ux43aux446ux438ux438-ux44dux440ux43cux438ux442ux43eux432ux43eux441ux442ux44c-ux438-ux43fux43eux442ux43eux43aux438}

Прежде чем говорить о статистике, конструкция обязана быть точной. Тест
15 проверяет на решётке \(16\times16\) с двумя вихрями \(q=\pm1\):
эрмитовость \(\|H-H^\dagger\|/\|H\| = 0\) (машинный нуль); класс
симметрии --- \(\tau(H)=\|H-H^*\|/\|H\|=0\), то есть в дираковской
калибровке матрица вещественная (класс GOE, как и предсказывает Дайсон);
поток через ячейку каждого вихря и через контрольную пустую ячейку --- с
точностью до \(10^{-14}\) рад. Тест 24 повторяет потоковую проверку для
одиночного вихря во всех 225 ячейках: 1/1 вихревых OK, 224/224 пустых
OK, максимум модуля потока на пустых ячейках \(1{,}47\cdot10^{-14}\).
Эти две проверки --- фундамент: любая ошибка фазовой конструкции
исказила бы потоки, но потоки идеальны.

\subsection{4.2 Байерс--Янг: целые вихри невидимы, дробные ---
физичны}\label{ux431ux430ux439ux435ux440ux441ux44fux43dux433-ux446ux435ux43bux44bux435-ux432ux438ux445ux440ux438-ux43dux435ux432ux438ux434ux438ux43cux44b-ux434ux440ux43eux431ux43dux44bux435-ux444ux438ux437ux438ux447ux43dux44b}

\begin{figure}
\centering
\includegraphics{../figures/fig14_byers_yang.png}
\caption{\textbf{Рис. 8.} Тест Байерса--Янга: целый заряд спектрально
невидим (\(3{,}5\cdot10^{-15}\)), дробный --- физичен (\(0{,}34\));
логарифмическая шкала.}
\end{figure}

\textbf{Теорема Байерса--Янга} (1961) в решёточной формулировке гласит:
спектр гамильтониана периодичен по добавке целого кванта потока на
ячейку --- целозарядный вихрь можно «откалибровать» (убрать) унитарным
преобразованием узлов, изменившим только фазы. Тест 25 измеряет это
прямо: две случайные конфигурации из четырёх нейтральных вихрей
\(q=\pm1\) на \(16\times16\) с открытыми границами при \(\alpha=0\) дают
максимальный спектральный сдвиг \(3{,}49\cdot10^{-15}\) относительно
чистой решётки --- целые вихри спектрально невидимы, тождество точное.
Одна дробная конфигурация \(q=0{,}3\) даёт сдвиг \(0{,}342\) --- дробный
вихрь несёт физический поток
\(4\pi\cdot0{,}3 = 1{,}2\pi \not\equiv 0 \pmod{2\pi}\) и реально
возмущает спектр. Следствие, которое определяет всю дальнейшую физику: в
дираковской калибровке GUE-статистику могут порождать только
\textbf{дробные} заряды либо фон \(\alpha\neq0\); в гладкой калибровке
\texttt{monumental} фазы комплексны при любом \(q\), и целые заряды
работают (раздел 5). Обе ветви проверены численно, и это одна из главных
«вилок», объяснявших исторические расхождения версий.

\subsection{4.3 Самодуальность Конна и хиральная симметрия
AIII}\label{ux441ux430ux43cux43eux434ux443ux430ux43bux44cux43dux43eux441ux442ux44c-ux43aux43eux43dux43dux430-ux438-ux445ux438ux440ux430ux43bux44cux43dux430ux44f-ux441ux438ux43cux43cux435ux442ux440ux438ux44f-aiii}

При \(\alpha=1/2\) гамильтониан Хофштадтера обладает
\textbf{самодуальностью Конна}: замена \(\alpha\to1/\alpha\)
эквивалентна преобразованию Фурье по узлам, и точка \(\alpha=1/2\)
неподвижна. На торе тест 17 измеряет: число нулевых мод при
\(\alpha=1/2\) равно ровно 4 (для чётных \(L\) это удвоенное число рёбер
неподвижной решётки --- аналитическое значение), спектральная
E→−E-симметрия --- дефект \(1{,}4\cdot10^{-15}\). \textbf{Хиральная
симметрия подрешёток} (тест 27): на двудольной квадратной решётке
оператор \(\Gamma\) (знак шахматной доски) удовлетворяет
\(\Gamma H\Gamma = -H\) для любого \(\alpha\) при отсутствии вихрей;
измеренный дефект при \(W=0\) равен \(0\) (точно), при \(W=4\) ---
\(5{,}9\cdot10^{-3}\) (мал, но ненулевой, поскольку беспорядок на узлах
нарушает двудольность), при двух вихрях без беспорядка --- снова ровно
\(0\): фазовые вихри двудольность сохраняют. Дополнительный
аналитический факт теста 17: при \(\alpha=1/3\) спектр имеет щель на
\(E=0\) (нулевых мод 0), а изоспектральность
\(\alpha\leftrightarrow1/\alpha\) на конечном торе нарушается (дефект
\(0{,}30\)) --- она точна только в блочном (блоховском) смысле; это
честно документировано, так как конечномерное нарушение не противоречит
бесконечномерной теореме.

\subsection{4.4 Число Черна: C₁ =
2}\label{ux447ux438ux441ux43bux43e-ux447ux435ux440ux43dux430-cux2081-2}

Тест 20 вычисляет первым зонным \textbf{число Черна} методом
Фукуи--Хацугаи--Сузуки на торе \(6\times6\) при \(\alpha=1/2\) (без
вихрей): сетка \(8\times8\) точек скрученных граничных условий, сечение
Берри нижней зоны даёт \(C_1 = 2\). Монография предсказывала \(|C_1|=1\)
на подзону; при \(\alpha=1/2\) зоны Хофштадтера сливаются попарно, и
«нижняя зона» содержит две подзоны с \(C_1=\pm1\) одинакового знака
вклада --- итог \(|C_1|=2\) согласован с предсказанием при правильном
учёте вырождения. Существенно одно: \(C_1\neq0\) --- фаза реализует
квантовый эффект Холла, топология нетривиальна, и защищённые краевые
состояния обязаны существовать (на открытых границах они видны как
«ножки» бабочки, рис. 20).

\subsection{4.5 Тор и нарушение
T-симметрии}\label{ux442ux43eux440-ux438-ux43dux430ux440ux443ux448ux435ux43dux438ux435-t-ux441ux438ux43cux43cux435ux442ux440ux438ux438}

Тест 26 объединяет потоковые и статистические проверки на периодической
решётке: тор \(16\times16\), одиночный вихрь \(q=1\) с
калибровочно-совместным условием \(q(2-N_x) = -14\in\mathbb{Z}\),
периодическая дираковская нить вдоль +x. Потоки: вихревая ячейка и все
255 пустых --- OK; эрмитовость --- OK; в гладкой калибровке для
статистики \(\tau(H) = 0{,}2252\) --- T-симметрия нарушена (класс GUE),
и \(\langle r\rangle = 0{,}5857\) --- в полосе GUE
\(0{,}5992\pm0{,}022\)-калибровки. Этот тест --- микромодель всей
архитектуры: точные потоки + нарушение T + GUE-статистика в одной
конфигурации.

\section{5. Блок 3: GUE-классификация AB-облака (тесты 16, 26, 29, 33,
34)}\label{ux431ux43bux43eux43a-3-gue-ux43aux43bux430ux441ux441ux438ux444ux438ux43aux430ux446ux438ux44f-ab-ux43eux431ux43bux430ux43aux430-ux442ux435ux441ux442ux44b-16-26-29-33-34}

\subsection{5.1 Роль калибровки: почему одна и та же модель даёт GOE и
GUE}\label{ux440ux43eux43bux44c-ux43aux430ux43bux438ux431ux440ux43eux432ux43aux438-ux43fux43eux447ux435ux43cux443-ux43eux434ux43dux430-ux438-ux442ux430-ux436ux435-ux43cux43eux434ux435ux43bux44c-ux434ux430ux451ux442-goe-ux438-gue}

Тест 16 (модель Монтгомери) в дираковской калибровке при \(N_v=2\),
\(q=\pm1\), \(\alpha=1/2\) даёт
\(\langle r\rangle = 0{,}5133\pm0{,}0353\) --- между GOE (0,531) и GUE
(0,600), с монте-карловским p-значением \(0{,}00\) против GUE-ансамбля
из 30 матриц \(256\times256\). Причина ясна из 4.2: целые заряды в этой
калибровке калибровочно невидимы, матрица почти вещественная, и
«потолок» GOE не может быть преодолён никаким беспорядком --- это
матричная, а не шумовая T-симметрия. Тест 16 повторён в гладкой
калибровке внутри теста 33: при тех же позициях вихрей
\(\langle r\rangle\) поднимается до \(0{,}585\). Мораль для
воспроизводимости: \textbf{калибровка фаз --- часть физической
спецификации}; сравнение версий сюиты без указания калибровки
бессмысленно, и в заголовке каждого отчёта она теперь указана.

\subsection{5.2 Многореализационный бутстрап (тест
33)}\label{ux43cux43dux43eux433ux43eux440ux435ux430ux43bux438ux437ux430ux446ux438ux43eux43dux43dux44bux439-ux431ux443ux442ux441ux442ux440ux430ux43f-ux442ux435ux441ux442-33}

\begin{figure}
\centering
\includegraphics{../figures/fig11_r_bootstrap.png}
\caption{\textbf{Рис. 9.} Тест 33: три реализации \(\langle r\rangle\)
для \(q=1\) (тор) и \(q=0{,}3\) (открытые границы) на фоне GUE- и
пуассоновской линий.}
\end{figure}

Основной статистический результат работы. Конфигурация: тор
\(16\times16\), два вихря \(q=\pm1\), \(\alpha=1/2\), беспорядок
\(W=4\), гладкая калибровка, три независимые реализации (случайные
позиции вихрей, разные зерна беспорядка). Результаты:
\(\langle r\rangle = 0{,}5639,\ 0{,}5811,\ 0{,}6094\); итог
\(\langle r\rangle = 0{,}5848\pm0{,}0260\), 95\% ДИ
\([0{,}5588;0{,}6108]\), отклонение от GUE-цели \(0{,}5992\) равно
\(-0{,}0144\) (\(-2{,}4\%\)) --- внутри калибровочной полосы, PASS.
Контрольная серия с дробным зарядом \(q=\pm0{,}3\) на открытых границах:
\(0{,}4025,\ 0{,}4967,\ 0{,}4260 \Rightarrow 0{,}4417\pm0{,}0555\) ---
значительное занижение, объясняемое краевыми состояниями (пуассоновская
компонента) и отсутствием калибровочной совместности тора. Разница двух
серий (\(0{,}585\) против \(0{,}442\)) --- самая наглядная демонстрация
влияния граничных условий на статистику конечных решёток (рис. 11).

\subsection{5.3 Скольжение по размерам (тест
34)}\label{ux441ux43aux43eux43bux44cux436ux435ux43dux438ux435-ux43fux43e-ux440ux430ux437ux43cux435ux440ux430ux43c-ux442ux435ux441ux442-34}

\begin{figure}
\centering
\includegraphics{../figures/fig10_r_L_scaling.png}
\caption{\textbf{Рис. 10.} Тест 34: сходимость \(\langle r\rangle\) с
размером решётки; обе ветки выходят на GUE-плато \(\approx0{,}60\).}
\end{figure}

Скейлинг \(\langle r\rangle(L)\) для \(L = 10, 20, 30, 50\) (две
реализации на точку, беспорядок \(W=4\)): дробная ветка \(q=0{,}3\)
(открытые границы) ---
\(0{,}5482 \to 0{,}5921 \to 0{,}6048 \to 0{,}5953\); целая ветка \(q=1\)
(тор) --- \(0{,}5739 \to 0{,}5681 \to 0{,}6042 \to 0{,}5864\). Обе ветки
при \(L\ge30\) выходят на плато \(0{,}59\)--\(0{,}605\) в пределах двух
стандартных ошибок от GUE \(0{,}5996\); дробная ветка монотонна и её
плато совпадает с предсказанием монографии
(\(0{,}487\to0{,}562\to0{,}595\to0{,}594\)) по форме сходимости ---
значение при \(L=10\) у нас выше (краевые эффекты при малом \(L\) у нас
слабее благодаря конкретной конфигурации вихрей), далее тренды
совпадают. Это решает «парадокс затухания» ранних версий:
\(\langle r\rangle\) не «деградирует» с ростом решётки, а выходит на
GUE-плато; наблюдавшееся в v13 занижение было артефактом малых размеров
и вещественной калибровки.

\subsection{5.4 Сводная мера f\_GUE (тест
29)}\label{ux441ux432ux43eux434ux43dux430ux44f-ux43cux435ux440ux430-f_gue-ux442ux435ux441ux442-29}

Тест 29 сводит статистику к одной шкале: при \(16\times16\), \(N_v=2\),
\(q=\pm1\), \(\alpha=1/2\), \(W=4\) значение
\(\Sigma^2(L{=}2) = 0{,}8377\) даёт \(f_{\rm GUE} = 0{,}648\) при пороге
PASS \(0{,}90\) --- WARN. Это не противоречит тесту 33: \(\Sigma^2\) при
малом \(L\) чувствительна к «жёсткому» ядру спектра и краям, где смесь
ещё не универсальна; порог \(0{,}90\) был откалиброван для ансамблей
\(\ge100\times100\). В блоке 5 мы увидим ту же картину в терминах
\(K(t)\): малые ансамбли дают правильную \emph{тенденцию}, но не
проходят жёсткие пороги, откалиброванные на больших. Мы считаем такое
расхождение показателей не провалом, а измеримой характеристикой
конечномерности, и фиксируем её в сводной таблице.

\section{6. Блок 4: прямое сравнение с нулями ζ (тесты 28, 35,
36)}\label{ux431ux43bux43eux43a-4-ux43fux440ux44fux43cux43eux435-ux441ux440ux430ux432ux43dux435ux43dux438ux435-ux441-ux43dux443ux43bux44fux43cux438-ux3b6-ux442ux435ux441ux442ux44b-28-35-36}

\subsection{6.1 Двухпроходное сравнение нулевых выборок (тест
28)}\label{ux434ux432ux443ux445ux43fux440ux43eux445ux43eux434ux43dux43eux435-ux441ux440ux430ux432ux43dux435ux43dux438ux435-ux43dux443ux43bux435ux432ux44bux445-ux432ux44bux431ux43eux440ux43eux43a-ux442ux435ux441ux442-28}

\begin{figure}
\centering
\includegraphics{../figures/fig15_berry_R2_cutoff.png}
\caption{\textbf{Рис. 11.} Берри-коррекция \(R_2(0;T)\):
конечновыборочный обрез и медленное приближение к асимптоте \(-1\);
вертикали --- высоты двух проходов (1000/5000 нулей).}
\end{figure}

Тест 28 --- тест Берри-поправки к \(R_2(0)\) --- в v18/v19 выполняется в
двух прогонах: первые 1000 нулей (\(T_{\max}\approx1419\)) и первые 5000
нулей (\(T_{\max}\approx5448\)). Сумма по простым
\(\sum_p\sum_{k\le5}p^{-k/2}/(k\ln p) = 4{,}2079\) даёт корректированное
теоретическое значение \(R_2(0) = -2{,}3394\); ожидаемая
конечновыборочная добавка \(\sim T^{-1/2}\) падает от \(0{,}0019\)
(\(T=1419\)) до \(0{,}0009\) (\(T=5448\)). Тест проверяет точное
воспроизведение формулы в обоих прогонах (PASS, машинная точность) и
демонстрирует сам масштаб берри-эффекта: даже корректированное
предсказание \(R_2(0)\) на достижимых высотах далеко от «школьного»
значения \(-1{,}007\) асимптотики. Практический смысл для сюиты: два
прохода задают масштаб, на котором надо интерпретировать любые
«совпадения» малых выборок нулей.

\subsection{6.2 Тест Монтгомери: AB-облако против ζ-5000 (тест
35)}\label{ux442ux435ux441ux442-ux43cux43eux43dux442ux433ux43eux43cux435ux440ux438-ab-ux43eux431ux43bux430ux43aux43e-ux43fux440ux43eux442ux438ux432-ux3b6-5000-ux442ux435ux441ux442-35}

\begin{figure}
\centering
\includegraphics{../figures/fig08_r2_pair.png}
\caption{\textbf{Рис. 12.} Парная корреляция \(R_2(s)\): AB-облако
против нулей \(\zeta\), GUE и Пуассона --- корреляционная дыра
воспроизводится обеими ветвями.}
\end{figure}

Прямое лобовое сравнение: развёрнутые спейсинги AB-облака (2 реализации,
304 зазора) против развёрнутых спейсингов первых 5000 нулей ζ (4999
зазоров, симметричная скользящая развёртка). Результаты: двухвыборочный
KS \(D = 0{,}1173\), \(p = 6{,}7\cdot10^{-4}\) --- гипотеза полного
совпадения распределений отвергнута; средний модуль разности парных
корреляций \(\langle|R_2^{\rm AB}(s)-R_2^{\zeta}(s)|\rangle = 0{,}147\)
при пороге \(0{,}10\). Но та же таблица \(R_2(s)\) (29 значений \(s\) от
0,05 до 2,85; рис. 8) показывает главное: корреляционная дыра есть у
обоих --- \(R_2^{\rm AB}(0{,}05) = 0{,}00\),
\(R_2^{\zeta}(0{,}05) = 0{,}005\) при пуассоновской единице; обе кривые
идут параллельно GUE-предсказанию и обе проходят проверку «ближе к GUE,
чем к Пуассону» (\(d_{\rm GUE} = 0{,}140 < d_{\rm Pois} = 0{,}227\)).
Честная формулировка итога: \textbf{AB-облако воспроизводит качественную
и полуколичественную структуру корреляций нулей (дыру, излом при
\(s\approx1\), выход на плато), но не совпадает с ними поточечно} ---
304 зазора не хватает ни для устойчивой \(R_2\), ни для KS-совпадения.
Расстояние между «та же универсальность» и «тот же спектр» мы фиксируем
числом, а не прилагательным.

\subsection{6.3 Формфактор K(t) (тест
36)}\label{ux444ux43eux440ux43cux444ux430ux43aux442ux43eux440-kt-ux442ux435ux441ux442-36}

\begin{figure}
\centering
\includegraphics{../figures/fig09_K_form_factor.png}
\caption{\textbf{Рис. 13.} Спектральный формфактор \(K(t)\): рампа+плато
GUE, данные AB-облака и нулей \(\zeta\); малые ансамбли дают правильную
тенденцию с выбросами.}
\end{figure}

Спектральный формфактор --- самая требовательная диагностика: для
устойчивой оценки нужны сотни независимых зазоров и точная развёртка.
Наши 306 собственных значений дают RMS-отклонение от GUE-рампы
\(0{,}93\) (порог PASS \(0{,}30\)) и корреляцию с теоретической кривой
\(0{,}018\) (порог \(0{,}50\)) --- WARN по обеим метрикам; форма кривой
(рис. 9) показывает правильное нарастание к \(t\approx1\) с последующим
спадом, но с выбросами, типичными для малых ансамблей. Кривая нулей ζ на
том же графике (по 5000 нулей) имеет RMS до GUE \(0{,}86\) --- то есть
даже эталонные нули на этих масштабах не проходят жёсткий порог:
диагностика информативна, но калибрована под ансамбли на порядки больше.
Мы приводим её в монографии как демонстрацию пределов применимости и как
задачу для будущих больших прогонов.

\section{7. Блок 5: дираковская динамика и неэрмитовость (тесты 19,
30--32,
37)}\label{ux431ux43bux43eux43a-5-ux434ux438ux440ux430ux43aux43eux432ux441ux43aux430ux44f-ux434ux438ux43dux430ux43cux438ux43aux430-ux438-ux43dux435ux44dux440ux43cux438ux442ux43eux432ux43eux441ux442ux44c-ux442ux435ux441ux442ux44b-19-3032-37}

\subsection{7.1 Дираковский конус (тесты 19,
31)}\label{ux434ux438ux440ux430ux43aux43eux432ux441ux43aux438ux439-ux43aux43eux43dux443ux441-ux442ux435ux441ux442ux44b-19-31}

\begin{figure}
\centering
\includegraphics{../figures/fig12_dirac_cone.png}
\caption{\textbf{Рис. 14.} Дираковский конус:
\(E_{\min} = 0{,}02 + 11{,}83/L\), \(R^2 = 0{,}9997\) --- линейность по
\(1/L\) при четырёх нулевых модах во всех точках.}
\end{figure}

Вихрь в AB-облаке при \(\alpha=1/2\) ведёт себя как релятивистский
фермион на прямой: спектр вблизи нуля линеен. Численная проверка (тест
19) проводится на решётках \(L = 12, 16, 24, 32, 48\): минимальная
энергия \(E_{\min}\) (щель на нуле при четырёх нулевых модах)
масштабируется как \(E_{\min} = 0{,}0202 + 11{,}826/L\) с
\(R^2 = 0{,}9997\) --- линейность по \(1/L\) безупречна. Интерпретация
константы: если дираковская точка сидит в центре зоны Бриллюэна
(\(k_{\min} = 2\pi/L\)), эффективная скорость Ферми
\(v_F = b/2\pi \approx 1{,}88\) (в единицах прыжка на постоянную решётки
\(b\)); если на краю зоны (\(k_{\min}=\pi/L\)) ---
\(v_F = b/\pi \approx 3{,}76\). Историческое значение \(\sim0{,}125\)
относилось к другой нормировке решётки; в текущей нормировке сюиты
фиксируем именно эти два числа и указываем, что выбор между ними ---
вопрос идентификации точки контакта зон, на которую указывает фазовая
структура нулевых мод. Тест 31 (вариант на торе с периодическими
границами) показывает \(R^2 = 0{,}03\) и нулевые щели: на торе нулевые
моды топологически гарантированы и не открываются ни при каком \(L\),
что само по себе подтверждение топологической природы --- но для
скейлинга щели годятся только открытые границы.

\subsection{7.2 Дираковский провал плотности состояний (тест
30)}\label{ux434ux438ux440ux430ux43aux43eux432ux441ux43aux438ux439-ux43fux440ux43eux432ux430ux43b-ux43fux43bux43eux442ux43dux43eux441ux442ux438-ux441ux43eux441ux442ux43eux44fux43dux438ux439-ux442ux435ux441ux442-30}

\begin{figure}
\centering
\includegraphics{../figures/fig13_dirac_dip.png}
\caption{\textbf{Рис. 15.} Дираковский провал DOS при \(\alpha=1/2\):
контраст \(20\times\) в обе стороны (логарифмическая шкала).}
\end{figure}

В точке \(\alpha=1/2\) спектр Хофштадтера имеет аналитическую
особенность: плотность состояний вблизи \(E=0\) обращается в ноль
линейно (\(\rho(E)\sim|E|\) дираковской точке), тогда как при
\(\alpha\neq1/2\) спектр типичен. Тест 30 измеряет плотность состояний в
центральной полосе \(|E|<0{,}5\) на решётке \(20\times20\) без
беспорядка: \(\rho(\alpha=0{,}4) = 0{,}200\),
\(\rho(\alpha=0{,}5) = 0{,}010\), \(\rho(\alpha=0{,}6) = 0{,}200\) ---
контраст ровно \(20\times\) в обе стороны. Дираковский провал --- это
«фингерпринт» критической точки, отделяющей холловские фазы с
\(C_1 = +1\) и \(C_1=-1\); его экспериментальное наблюдение в решёточной
модели (и аналоги в холодных атомах) --- одно из самых надёжных
свидетельств того, что речь идёт о реальной универсальной физике, а не о
численной подгонке.

\subsection{\texorpdfstring{7.3 Хатано--Нельсон: уход от \(\sigma=1/2\)
(тест
32)}{7.3 Хатано--Нельсон: уход от \textbackslash sigma=1/2 (тест 32)}}\label{ux445ux430ux442ux430ux43dux43eux43dux435ux43bux44cux441ux43eux43d-ux443ux445ux43eux434-ux43eux442-sigma12-ux442ux435ux441ux442-32}

\begin{figure}
\centering
\includegraphics{../figures/fig16_hatano_nelson.png}
\caption{\textbf{Рис. 16.} Хатано--Нельсон: \(\sigma=0{,}5\) ---
вещественный спектр; \(\sigma=0{,}7\) --- весь спектр в комплексной
плоскости (скин-эффект).}
\end{figure}

Параметр \(\sigma\) в обобщённой модели задаёт «дробную» часть фазы
прыжков; при \(\sigma=1/2\) модель эрмитова. Тест 32 сравнивает
\(\sigma=0{,}5\) и \(\sigma=0{,}7\) на \(16\times16\): при
\(\sigma=0{,}5\) максимум \(|\mathrm{Im}\,E| = 0\) (спектр строго
вещественный); при \(\sigma=0{,}7\) весь спектр из 256 уровней
комплексный, максимум \(|\mathrm{Im}\,E| = 0{,}623\), параметр
скин-эффекта Хатано--Нельсона \(g = (\sigma-1/2)\ln|\gamma| = 0{,}322\).
Смысл: \textbf{уход от \(\sigma=1/2\) переводит систему из унитарного
(GUE) класса в диссипативный (некоррелированный) класс; локализация
собственных векторов на границе (скин-эффект) разрушает объёмную
статистику, \(\langle r\rangle\) на неэрмитовой ветке обнуляется (все
зазоры «парные» в комплексной плоскости)}. Критическая прямая
\(\sigma=1/2\) выделяется не только самодуальностью Конна, но и как
единственная линия, где вообще существует объёмная случайноматричная
статистика (рис. 13, 16). Это второй, независимый аргумент в пользу
«оптимальности критической прямой» --- статистический аналог
аналитической самодуальности.

\subsection{7.4 Байтовая устойчивость (тест
37)}\label{ux431ux430ux439ux442ux43eux432ux430ux44f-ux443ux441ux442ux43eux439ux447ux438ux432ux43eux441ux442ux44c-ux442ux435ux441ux442-37}

Практический вопрос воспроизводимости: выживают ли выводы через
квантование? Тест 37 проверяет, что \(\langle r\rangle\) устойчив к
квантованию собственных значений на 256 уровней (изменение \(0{,}0015\)
при допуске \(0{,}02\)) и что байтовый поток генератора имеет энтропию
Шеннона \(7{,}9991\) бит при максимуме 8 (\(\chi^2\)-критерий через ЦПТ:
\(z=-0{,}234\), \(|z|<3\)). Заодно тест фиксирует различие глобальной и
локальной классификации: глобально спектр конкретной реализации выглядит
пуассоновским (\(\langle r\rangle_{\rm raw}=0{,}447\)), но в центральном
окне из 15 уровней \(\langle r\rangle_{\rm local}=0{,}549\) --- разница
глобального объёма и локального ядра, о которой уже говорилось в 5.4.

\section{8. Блок 6: аналитические интерпретации (тесты 21--23 и
обобщения)}\label{ux431ux43bux43eux43a-6-ux430ux43dux430ux43bux438ux442ux438ux447ux435ux441ux43aux438ux435-ux438ux43dux442ux435ux440ux43fux440ux435ux442ux430ux446ux438ux438-ux442ux435ux441ux442ux44b-2123-ux438-ux43eux431ux43eux431ux449ux435ux43dux438ux44f}

Этот блок --- мост от измерений к аналитическим структурам. Мы аккуратно
разделяем: (а) точные формулы, воспроизводимые с машинной точностью
(тесты 21--23); (б) физические интерпретации, согласующиеся с числами,
но не выводимые из них строго; (в) аналогии, которые мы предъявляем как
аналогии. Такая трёхуровневая подача --- принципиальная часть
методологии.

\subsection{\texorpdfstring{8.1 Спинорная фаза \(\gamma^{\ast}\): 90°
через мнимую единицу (тест
21)}{8.1 Спинорная фаза \textbackslash gamma\^{}\{\textbackslash ast\}: 90° через мнимую единицу (тест 21)}}\label{ux441ux43fux438ux43dux43eux440ux43dux430ux44f-ux444ux430ux437ux430-gammaast-90-ux447ux435ux440ux435ux437-ux43cux43dux438ux43cux443ux44e-ux435ux434ux438ux43dux438ux446ux443-ux442ux435ux441ux442-21}

\begin{figure}
\centering
\includegraphics{../figures/fig18_gamma_phase.png}
\caption{\textbf{Рис. 17.} Спинорная фаза
\(\gamma^{\ast} = a_C + i b_C\): аргумент \(89{,}87°\) против эталонных
\(90°\); доминирование \(b_C\) прижимает фазу к \(\pi/2\).}
\end{figure}

Определим \(\delta_C = \pi/7 = 0{,}448799\) рад --- см. раздел 8.2 --- и
\(b_2(K3) = 22\) (вторая число Бетти K3-поверхности). Тест 21 проверяет
две вещественные величины: \textbf{вещественный тормозящий член}
\(a_C(\gamma^{\ast}) = \delta_C^5/b_2(K3) = 8{,}276\cdot10^{-4}\) и
\textbf{мнимый берри-член}
\(b_C(\gamma^{\ast}) = 1-\cos(2\delta_C) = 2\sin^2\delta_C = 0{,}376510\);
обе воспроизводятся с машинной точностью. Комплексная спинорная фаза
\(\gamma^{\ast} = a_C + i b_C\) имеет модуль
\(|\gamma^{\ast}| = 0{,}3765\) и аргумент \(89{,}874°\) --- отклонение
от \(\pi/2\) всего \(0{,}126°\) (\(2{,}2\cdot10^{-3}\) рад). Структура
здесь та же, что в тесте 21: так как \(b_C \gg a_C\), аргумент
\(\gamma^{\ast}\) автоматически прижимается к \(90°\) --- «переход через
90° реализуется мнимой единицей»: фаза как геометрический объект
сконструирована так, что её релятивистская (спинорная) природа
выражается чисто мнимым доминированием. Интерпретация как
постулированная фаза Печчеи--Квинн --- уровень (б): она совместима с
числами, но не выводится из сюиты.

\subsection{\texorpdfstring{8.2 АБ-фаза \(\Phi_{AB} = \pi/7\) и дробный
заряд (тест
22)}{8.2 АБ-фаза \textbackslash Phi\_\{AB\} = \textbackslash pi/7 и дробный заряд (тест 22)}}\label{ux430ux431-ux444ux430ux437ux430-phi_ab-pi7-ux438-ux434ux440ux43eux431ux43dux44bux439-ux437ux430ux440ux44fux434-ux442ux435ux441ux442-22}

Точная формула: для частицы с зарядом \(q/e = 1/14\), огибающей трубку
одного кванта потока, фаза Ааронова--Бома равна
\(\Phi_{AB} = 2\pi/14 = \pi/7 = 0{,}448799\) рад \(= 25{,}714°\). Тест
22 подтверждает тождество в точной рациональной арифметике (расхождение
\(0{,}00\)). Величина \(\delta_C=\pi/7\) --- это и есть \(\Phi_{AB}\),
то есть \(q/e=1/14\) --- постулат дробного заряда, вокруг которого
строится вся аналитическая надстройка. Существенно, что \(\pi/7\)
связывается с гауссовой суммой \(b_7=(-1+i\sqrt7)/2\) таблицы характеров
PSL(2,7) (раздел 8.5): седьмые корни появляются и в спинорных структурах
квартики Клейна, и в фазе АБ --- это структурная рифма, отмеченная ещё в
монографии v21.

\subsection{8.3 Комплексный фрактальный фактор CF и логарифмическая
параболизация (тест
23)}\label{ux43aux43eux43cux43fux43bux435ux43aux441ux43dux44bux439-ux444ux440ux430ux43aux442ux430ux43bux44cux43dux44bux439-ux444ux430ux43aux442ux43eux440-cf-ux438-ux43bux43eux433ux430ux440ux438ux444ux43cux438ux447ux435ux441ux43aux430ux44f-ux43fux430ux440ux430ux431ux43eux43bux438ux437ux430ux446ux438ux44f-ux442ux435ux441ux442-23}

Для критического решения с дискретным масштабным отношением
\(\lambda_{\rm DSI} = b_2(K3) = 22\) тест 23 вычисляет
\textbf{комплексный фрактальный фактор}

\[CF = \exp\!\left(\frac{2\pi i}{\lambda_{\rm DSI}}\right)\cdot\exp\!\left(i\,\beta_{\rm imag}\,\ln\lambda_{\rm DSI}\right),\qquad \beta_{\rm imag} = 2 + \frac{1}{2\sin^2(\pi/7)} - 2\cos\frac{\pi}{7} = 2{,}854033,\]

что даёт \(CF = -0{,}9501 + 0{,}3119\,i\) с \(|CF| = 1\) по построению.
Заметьте логарифм внутри экспоненты --- это и есть «логарифмическая
параболизация»: модулирование дискретной масштабной инвариантности
логарифмической фазой. Отдельная честная оговорка теста: захардкоженное
в поздних tex-источниках значение
\(CF_{\rm hard} = 0{,}9786 + 0{,}0390\,i\) (модуль \(0{,}9794\), откуда
\(c_{AB} = 1-|CF_{\rm hard}| = 0{,}0206\)) --- это \textbf{эмпирический
вход}, аналог \(\alpha_s\) в КХД, а формула --- структурное
(топологическое) предсказание; амплитуды различаются именно на
\(c_{AB}\), и тест фиксирует оба числа с их статусами.

\subsection{8.4 Нарушение T-симметрии как происхождение GUE: аналогия с
чёрными
дырами}\label{ux43dux430ux440ux443ux448ux435ux43dux438ux435-t-ux441ux438ux43cux43cux435ux442ux440ux438ux438-ux43aux430ux43a-ux43fux440ux43eux438ux441ux445ux43eux436ux434ux435ux43dux438ux435-gue-ux430ux43dux430ux43bux43eux433ux438ux44f-ux441-ux447ux451ux440ux43dux44bux43cux438-ux434ux44bux440ux430ux43cux438}

Числа блоков 3--5 складываются в аналитическую картину Дайсона: GOE ↔
T-инвариантность, GUE ↔ её нарушение. AB-облако реализует обе ветви в
одной модели: дираковская калибровка при целых зарядах --- вещественная
матрица, \(\tau=0\), GOE-потолок (тесты 15, 16); гладкая калибровка или
дробные заряды --- \(\tau\approx0{,}225\), GUE-статистика (тесты 26,
33). Предложенная в обсуждениях гипотеза усиливает эту линию:
\textbf{нарушение T-симметрии --- та же причина, по которой статистика
чёрных дыр (где время искривлено и обращение не является симметрией
термодинамически) описывается унитарными ансамблями}. Мы подчёркиваем
статус этой конструкции: это не следствие сюиты, а интерпретационная
рамка, которую сюита не противоречит --- и которая делает предсказание
«любая фазовая модель с комплексными прыжками и без интегрируемости даст
\(\langle r\rangle\to0{,}6\)» проверяемым на независимых материалах.

\subsection{8.5 Керр-аналог, эргосфера и PSL(2,7) ⊂ Monster
(обзор)}\label{ux43aux435ux440ux440-ux430ux43dux430ux43bux43eux433-ux44dux440ux433ux43eux441ux444ux435ux440ux430-ux438-psl27-monster-ux43eux431ux437ux43eux440}

Две интерпретации из версии 21 мы сохраняем в обзорном статусе.
\textbf{Керр-аналог}: фазовое поле вихря задаёт «эргорегион» ---
область, где фазовая скорость относительно решётки меняет знак, как
эргосфера Керра; численная проверка требует динамики, а не статического
спектра, и остаётся задачей. \textbf{Монстролический мостик}: характер
\(\chi_2\) Монстра размерности 196883 при ограничении на подгруппу
PSL(2,7) (вложение через максимальную подгруппу
\(2^{3+6+12+18}.(L_3(2)\times3S_6)\)) приближённо разлагается как
\(1456\cdot\chi_1 + 3334\cdot\chi_3 + 3334\cdot\chi_3' + 6784\cdot\chi_6 + 8081\cdot\chi_7 + 9770\cdot\chi_8\)
(сумма размерностей 196891, отклонение 0,004\%); кратности трактовались
бы как правила отбора многочастичных вихревых состояний. Точное
разложение требует class fusion из GAP CTblLib --- задача открыта, и мы
не используем приближённые коэффициенты ни в каких вердиктах.

\section{9. Синтез: нули ζ как коды допустимых
состояний}\label{ux441ux438ux43dux442ux435ux437-ux43dux443ux43bux438-ux3b6-ux43aux430ux43a-ux43aux43eux434ux44b-ux434ux43eux43fux443ux441ux442ux438ux43cux44bux445-ux441ux43eux441ux442ux43eux44fux43dux438ux439}

\begin{figure}
\centering
\includegraphics{../figures/fig19_vortex_texture.png}
\caption{\textbf{Рис. 18.} Фазовая текстура вихревой пары \(q=\pm1\) на
фоне калибровки Ландау: «облако» фаз, порождающее наблюдаемую
статистику.}
\end{figure}

\begin{figure}
\centering
\includegraphics{../figures/fig20_hofstadter_butterfly.png}
\caption{\textbf{Рис. 19.} Бабочка Хофштадтера \(E(\alpha)\); вертикаль
--- самодвойственная точка \(\alpha=1/2\), в которой сосредоточены
ключевые результаты монографии.}
\end{figure}

Соберём цепочку. (1) Спектр AB-облака при \(\alpha=1/2\) в гладкой
калибровке имеет GUE-статистику, сходимую с размером:
\(\langle r\rangle: 0{,}548\to0{,}595\) при \(L: 10\to50\) (тесты 33,
34). (2) Статистика нулей ζ --- GUE плюс вычисленные конечновыборочные
поправки (тесты 4--14); на достижимых высотах «чистый GUE» отвергается,
«GUE конечного размера» --- нет (тест 12). (3) Прямое сравнение
показывает совпадение качественной структуры \(R_2\) (тест 35) ---
корреляционная дыра Монтгомери воспроизводится вихревым облаком. (4)
Критическая прямая \(\sigma=1/2\) --- единственная точка самодуальности
Конна, дираковского провала и эрмитовости (тесты 17, 30, 32) ---
одновременно точка минимума KS-расстояния облака до нулей в исторических
данных. (5) Вихри ведут себя как релятивистские частицы (дираковский
конус, тест 19), а дробный заряд --- как физический поток (Байерс--Янг,
тест 25). Итоговый нарратив, в котором «нули ζ суть коды допустимых
энергетических состояний фазового резонатора», получает статус
\textbf{непротиворечивой, количественно проверенной программы}, а не
доказанной теоремы: мы знаем, какие звенья машинно точны (потоки,
тождества), какие статистически устойчивы (скейлинг, бутстрап) и какие
остаются открытыми (асимптотика \(T\to\infty\), точечное совпадение
\(R_2\)).

\section{10. Методология верификации: что именно гарантирует
«PASS»}\label{ux43cux435ux442ux43eux434ux43eux43bux43eux433ux438ux44f-ux432ux435ux440ux438ux444ux438ux43aux430ux446ux438ux438-ux447ux442ux43e-ux438ux43cux435ux43dux43dux43e-ux433ux430ux440ux430ux43dux442ux438ux440ux443ux435ux442-pass}

Версия 19 сюиты --- первый релиз, в котором движок отчётов архитектурно
изолирован от счётной части. История этого решения поучительна: в v18
баг в PDF-райтере (список строк был типизирован как
\texttt{String{[}{]}}, а в него помещались кортежи) вызывал
\texttt{MethodError} после завершения \textbf{каждого} теста; исключение
перехватывалось внешним обработчиком, вердикты обнулялись в WARN, а
второй проход двухпроходного режима не запускался вовсе --- «быстрая
проверка» деградировала до однопроходной, хотя внешне выглядела
работающей. Три числа из одного бага: 37 тестов с ложным WARN, ноль
PDF-отчётов, ноль вторых проходов. Исправление (FIX-R1\ldots R5 в
приложении A) сделало следующее: (1) PDF-райтер корректен; (2) отказ
любого отчётного формата не влияет на вердикт; (3) тест 28 (область
видимости переменной двухпроходного цикла) и тест 18 (калибровочная
совместимость тора при \(\alpha=1/3\) на решётках, кратных не шести)
исправлены; (4) быстрая проверка усечена до 5000 нулей и реально
выполняет оба прохода 16×16→32×32. Рефлексия: \textbf{вывод сюиты не
сильнее её weakest-инфраструктуры}, и «почти работающий» отчётный
конвейер опаснее неработающего, потому что создаёт видимость
верифицированности.

Практические гарантии текущей версии: каждый тест --- отдельный каталог
с полными логами; каждая цифра монографии трассируется к номеру теста;
вердикт воспроизводится командой вида
\texttt{julia\ ab\_cloud\_v19.jl\ -\/-test\ N\ -\/-no-two-pass};
двухпроходность включена по умолчанию для полного прогона и для
\texttt{-\/-quick}; генератор случайных чисел имеет фиксированные зерна
во всех тестах, где случайность нужна (воспроизводимость бит в бит).

\section{11. Выводы}\label{ux432ux44bux432ux43eux434ux44b}

\begin{enumerate}
\def\labelenumi{\arabic{enumi}.}
\tightlist
\item
  \textbf{Точность конструкции} (машинный уровень): эрмитовость 0;
  потоки дираковских нитей с точностью \(10^{-14}\); Байерс--Янг для
  целых зарядов \(3{,}5\cdot10^{-15}\); самодуальность Конна --- 4
  нулевые моды и E→−E с дефектом \(10^{-15}\); хиральность \(0{,}0\) при
  \(W=0\). Конструкция математически корректна --- это доказано, а не
  оценено.
\item
  \textbf{GUE-универсальность} (статистический уровень):
  \(\langle r\rangle = 0{,}5848\pm0{,}0260\) против GUE \(0{,}5992\)
  (отклонение \(-2{,}4\%\)); скейлинг \(L=10\to50\) монотонен и выходит
  на плато \(0{,}595\pm0{,}008\); спектр жёстче пуассоновской на всех
  масштабах \(L\) и ближе к GUE в 9/9 точках. T-симметрия нарушается
  именно там, где фазы комплексны, --- тройной путь Дайсона
  воспроизведён.
\item
  \textbf{Сравнение с нулями ζ} (честный уровень): качественная
  структура \(R_2(s)\) совпадает (дыра, плато; ближе к GUE, чем
  Пуассон), поточечного совпадения нет и не должно быть на ансамбле 304
  зазоров; конечновыборочные поправки Берри количественно объясняют все
  наблюдаемые отклонения от асимптотического GUE.
\item
  \textbf{Топология и динамика}: \(C_1=2\) (нетривиальный IQHE),
  дираковский конус \(R^2=0{,}9997\), провал DOS \(20\times\),
  скин-эффект при \(\sigma\neq1/2\); критическая прямая \(\sigma=1/2\)
  тройственно выделена (самодуальность + эрмитовость + провал).
\item
  \textbf{Аналитические интерпретации}: формулы \(\gamma^{\ast}\),
  \(\Phi_{AB}=\pi/7\), \(CF\) воспроизводятся с машинной точностью и
  образуют согласованную \(\pi/7\)-структуру; их физическая
  интерпретация (Печчеи--Квинн, монстрология) остаётся на уровне гипотез
  --- сюита не противоречит, но и не доказывает.
\item
  \textbf{Инфраструктура доверия}: v19 исправляет все обнаруженные
  дефекты отчётного конвейера; каждый вывод монографии трассируем к
  логам; две системы вердиктов (машинная и калибровочная) разделены
  явно.
\end{enumerate}

\section{12. Открытые
вопросы}\label{ux43eux442ux43aux440ux44bux442ux44bux435-ux432ux43eux43fux440ux43eux441ux44b}

\begin{enumerate}
\def\labelenumi{\arabic{enumi}.}
\tightlist
\item
  \textbf{Асимптотика \(b(N)\)}: степенной или логарифмический закон?
  Альтернативный фит лучше (\(R^2=0{,}9999\)), но диапазон \(N\)
  недостаточен; нужен прогон до \(N\sim10^6\) (данные Одлыжко 2M уже
  поддержаны загрузчиком).
\item
  \textbf{Точечное совпадение \(R_2\)}: ансамбль AB-облака нужно
  расширить до \(\ge10^4\) зазоров (большие решётки × многократные
  реализации × сдвиги вихрей) --- тогда двухвыборочный KS против ζ-5000
  станет информативным.
\item
  \textbf{Формфактор \(K(t)\)}: та же цель; пороги PASS
  (\(\mathrm{RMS}<0{,}30\), \(\mathrm{corr}>0{,}50\)) достигаются, по
  нашей оценке, при \(\ge3000\) зазоров.
\item
  \textbf{Интерфейс с 3D}: линии вихрей в стеке Хофштадтера (модуль
  3D-лаборатории) --- сохраняется ли GUE-плато по третьему измерению;
  есть предварительные данные, но не 37-тестовая верификация.
\item
  \textbf{Точное разложение 196883 по PSL(2,7)}: нужен class fusion
  через GAP CTblLib; до тех пор коэффициенты остаются приближением.
\item
  \textbf{Динамика вихрей и керр-аналог}: статический спектр не
  проверяет эргосферную интерпретацию; нужна симуляция Лагранжевой
  динамики точечных вихрей с АБ-фазами.
\end{enumerate}

\section{Приложение A. Код: архитектура и ключевые
модули}\label{ux43fux440ux438ux43bux43eux436ux435ux43dux438ux435-a.-ux43aux43eux434-ux430ux440ux445ux438ux442ux435ux43aux442ux443ux440ux430-ux438-ux43aux43bux44eux447ux435ux432ux44bux435-ux43cux43eux434ux443ux43bux438}

Полный код верификационной сюиты (файл \texttt{ab\_cloud\_v19.jl}, Julia
1.10, без внешних зависимостей --- включая собственные реализации
PNG-райтера с pHYs, векторного PDF, GIF89a с LZW, DOCX как stored-ZIP и
всей статистики) поставляется вместе с монографией в каталоге
\texttt{code/}; при запуске без аргументов открывает интерактивное меню,
ключ \texttt{f} --- быстрая проверка. Ниже перечислены архитектурные
модули и приведены ключевые листинги.

\textbf{Модули файла}: (1) конфигурация \texttt{Config} и двуязычные
сообщения; (2) загрузчики данных --- 50 000 встроенных нулей Одлыжко +
внешние файлы до 2M с мягким откатом; (3) статистическое ядро
(развёртка, KS, χ², AD, бутстрап, \(\Sigma^2\), \(\Delta_3\), \(R_2\),
\(K(t)\)); (4) строитель AB-облака
\texttt{build\_ab\_cloud\_hamiltonian} с тремя моделями фаз
(\texttt{:dirac}, \texttt{:monumental}) и проверкой потоков
\texttt{plaquette\_flux}; (5) 37 тестов \texttt{test*}; (6)
двухпроходный раннер \texttt{run\_test\_two\_pass}; (7) движок отчётов
v16/v19 (\texttt{rep\_*}, райтеры md/html/pdf/docx/svg/png/gif,
index.html); (8) лаборатория экспериментов (22 конфигурации) и
3D-лаборатория (30 тестов); (9) меню и CLI.

\textbf{Ключевая функция --- строитель гамильтониана} (сокращённо;
полный текст --- в поставляемом файле):

\begin{Shaded}
\begin{Highlighting}[]
\CommentTok{\# φ\_ij = 2π·α·j + Σ\_k q\_k·[arg(r\_i−r\_k) − arg(r\_j−r\_k)] + φ\^{}строка\_ij}
\CommentTok{\# связь (i→j) получает H\_ij = −exp(iφ\_ij); ячейка с вихрем k получает поток 2π·q\_k}
\KeywordTok{function} \FunctionTok{build\_ab\_cloud\_hamiltonian}\NormalTok{(cfg}\OperatorTok{::}\DataTypeTok{ABCloudConfig}\NormalTok{)}
\NormalTok{    N }\OperatorTok{=}\NormalTok{ cfg.Nx }\OperatorTok{*}\NormalTok{ cfg.Ny}
\NormalTok{    H }\OperatorTok{=} \FunctionTok{zeros}\NormalTok{(}\DataTypeTok{ComplexF64}\NormalTok{, N, N)}
    \ControlFlowTok{for}\NormalTok{ iy }\KeywordTok{in} \FloatTok{0}\OperatorTok{:}\NormalTok{(cfg.Ny}\OperatorTok{{-}}\FloatTok{1}\NormalTok{), ix }\KeywordTok{in} \FloatTok{0}\OperatorTok{:}\NormalTok{(cfg.Nx}\OperatorTok{{-}}\FloatTok{1}\NormalTok{)}
\NormalTok{        i }\OperatorTok{=}\NormalTok{ iy }\OperatorTok{*}\NormalTok{ cfg.Nx }\OperatorTok{+}\NormalTok{ ix}
        \ControlFlowTok{for}\NormalTok{ (dx, dy) }\KeywordTok{in}\NormalTok{ ((}\FloatTok{1}\NormalTok{,}\FloatTok{0}\NormalTok{),(}\FloatTok{0}\NormalTok{,}\FloatTok{1}\NormalTok{))            }\CommentTok{\# ближайшие соседи, обе ориентации}
\NormalTok{            jx, jy }\OperatorTok{=}\NormalTok{ (ix}\OperatorTok{+}\NormalTok{dx) }\OperatorTok{\%}\NormalTok{ cfg.Nx, (iy}\OperatorTok{+}\NormalTok{dy) }\OperatorTok{\%}\NormalTok{ cfg.Ny}
\NormalTok{            j }\OperatorTok{=}\NormalTok{ jy }\OperatorTok{*}\NormalTok{ cfg.Nx }\OperatorTok{+}\NormalTok{ jx}
            \ConstantTok{φ} \OperatorTok{=} \FloatTok{2}\NormalTok{π }\OperatorTok{*}\NormalTok{ cfg.alpha }\OperatorTok{*}\NormalTok{ jy              }\CommentTok{\# калибровка Ландау (по столбцу)}
            \ControlFlowTok{for}\NormalTok{ v }\KeywordTok{in}\NormalTok{ cfg.vortices                }\CommentTok{\# фазы вихрей (модель задаёт форму)}
                \ConstantTok{φ} \OperatorTok{+=} \FunctionTok{vortex\_phase}\NormalTok{(v, ix, iy, jx, jy, cfg.model)}
            \ControlFlowTok{end}
\NormalTok{            H[i}\OperatorTok{+}\FloatTok{1}\NormalTok{, j}\OperatorTok{+}\FloatTok{1}\NormalTok{] }\OperatorTok{{-}=} \FunctionTok{exp}\NormalTok{(}\ConstantTok{im} \OperatorTok{*} \ConstantTok{φ}\NormalTok{)}
\NormalTok{            H[j}\OperatorTok{+}\FloatTok{1}\NormalTok{, i}\OperatorTok{+}\FloatTok{1}\NormalTok{] }\OperatorTok{{-}=} \FunctionTok{exp}\NormalTok{(}\OperatorTok{{-}}\ConstantTok{im} \OperatorTok{*} \ConstantTok{φ}\NormalTok{)          }\CommentTok{\# эрмитово сопряжение}
        \ControlFlowTok{end}
\NormalTok{        H[i}\OperatorTok{+}\FloatTok{1}\NormalTok{, i}\OperatorTok{+}\FloatTok{1}\NormalTok{] }\OperatorTok{+=} \FunctionTok{disorder\_epsilon}\NormalTok{(ix, iy, cfg)   }\CommentTok{\# ε \textasciitilde{} U(−W/2, W/2), фиксированное зерно}
    \ControlFlowTok{end}
    \ControlFlowTok{return}\NormalTok{ H}
\KeywordTok{end}
\end{Highlighting}
\end{Shaded}

\textbf{Проверка потока через ячейку} (дираковская калибровка): сумма
фаз четырёх связей ячейки с вихрем даёт \(2\pi q \pmod{2\pi}\), пустой
--- 0; вычисляется через \texttt{plaquette\_flux(H,\ Nx,\ Ny,\ ix,\ iy)}
и проверяется во всех тестах 15/24/26.

\textbf{Вердикты}: \texttt{log\_result} пишет строку вердикта и в полный
журнал; \texttt{run\_test\_two\_pass} прогоняет тест на основной и
увеличенной решётке, восстанавливая конфигурацию в \texttt{finally}
(изоляция проходов); \texttt{rep\_finish!} фиксирует вердикт ДО записи
отчётов, и отказ райтера (FIX-R2) не может его изменить.

\textbf{Исправления v19 (полные комментарии в коде)}: FIX-R1 ---
типизация промежуточного списка строк в PDF-райтере (устраняет массовый
MethodError, обнулявший вердикты и второй проход); FIX-R2 --- try/catch
вокруг \texttt{\_write\_test\_report!}; FIX-R3/3b ---
\texttt{last\_R2}/\texttt{berry\_ref} в области видимости функции в
тесте 28; FIX-R4 --- скачок решётки к кратному 6 при α∈\{1/2,1/3\} в
тесте 18; FIX-R5/5b --- быстрая проверка: усечение ζ до 5000, отключение
GIF, CLI-флаг \texttt{-\/-quick}.

\section{Приложение B. Руководство по
воспроизведению}\label{ux43fux440ux438ux43bux43eux436ux435ux43dux438ux435-b.-ux440ux443ux43aux43eux432ux43eux434ux441ux442ux432ux43e-ux43fux43e-ux432ux43eux441ux43fux440ux43eux438ux437ux432ux435ux434ux435ux43dux438ux44e}

\begin{Shaded}
\begin{Highlighting}[]
\CommentTok{\# вся сюита (двухпроходная, \textasciitilde{}многие минуты):}
\ExtensionTok{julia}\NormalTok{ ab\_cloud\_v19.jl }\AttributeTok{{-}{-}test}\NormalTok{ all}
\CommentTok{\# отдельный тест:}
\ExtensionTok{julia}\NormalTok{ ab\_cloud\_v19.jl }\AttributeTok{{-}{-}test}\NormalTok{ 33 }\AttributeTok{{-}{-}no{-}two{-}pass}
\CommentTok{\# быстрая проверка 16×16 → 32×32 (ζ ≤ 5000, \textasciitilde{}минуты):}
\ExtensionTok{julia}\NormalTok{ ab\_cloud\_v19.jl }\AttributeTok{{-}{-}quick}
\CommentTok{\# двуязычность:}
\ExtensionTok{julia}\NormalTok{ ab\_cloud\_v19.jl }\AttributeTok{{-}{-}lang}\NormalTok{ ru }\AttributeTok{{-}{-}quick}
\CommentTok{\# отчёты: results/run\_\textless{}дата\textgreater{}/test\_NN\_\textless{}slug\textgreater{}/\{report.md,html,pdf,docx, plots/, logs/\}}
\end{Highlighting}
\end{Shaded}

Числа монографии соответствуют полному двухпроходному прогону
\textbf{v19 \texttt{run\_20260902\_134759}} от 2026-09-02 (37 тестов,
two-pass 72×72 → 96×96 с HARDCORE-аудитом второго прохода, серия pass-3
для тестов 19/29/31; Julia 1.12.0; 50 000 нулей Одлыжко). Итог первого
прохода: \textbf{32 PASS / 5 WARN} (тесты 4, 5, 6, 11, 12, 28, 35 ---
статистические WARN-вердикты «этажа сюрмиза Вигнера», калиброванные
двухвыборочными KS-строками против точного GUE-референса; семантика
сюиты не менялась); \textbf{все HARDCORE-проходы второго прохода ---
PASS} (суммарно 77 строк вердиктов). Полная поставка артефактов ---
отчёты (report.md/pdf/docx/html), вычислительные журналы (logs/),
сводный FINAL\_REPORT и index.html --- включена в репозиторий:
\texttt{results/run\_20260902\_134759/}. Дополнительно верификация всех
64 спинорных структур квартики Клейна выполнена эталонной реализацией
\texttt{verification/spinor64} (Приложение D). Ранний однопроходный
базлайн v18
(\texttt{results/verification\_run\_v18\_37tests\_2026-08-28.txt})
сохранён для истории.

\section{Приложение C. Указатель
графиков}\label{ux43fux440ux438ux43bux43eux436ux435ux43dux438ux435-c.-ux443ux43aux430ux437ux430ux442ux435ux43bux44c-ux433ux440ux430ux444ux438ux43aux43eux432}

Все графики --- 600 dpi, PNG, подписи на языке издания; файлы лежат в
\texttt{monographs/\textless{}lang\textgreater{}/figures/}. Рис. 1 ---
сходимость \(b(N)\) (два фита); рис. 2 --- гистограмма спейсингов ζ
против сюрмизов Вигнера; рис. 3 --- два фита затухания в
лог-координатах; рис. 4 --- бутстрап-распределение наклона с 95\% ДИ;
рис. 5 --- KS \(D(T_{\min})\) с критической линией; рис. 6 ---
\(\Sigma^2(L)\); рис. 7 --- \(\Delta_3(L)\); рис. 8 --- \(R_2(s)\):
AB-облако против ζ, GUE, Пуассона; рис. 9 --- формфактор \(K(t)\); рис.
10 --- скейлинг \(\langle r\rangle(L)\) для \(q=0{,}3\) и \(q=1\); рис.
11 --- бутстрап \(\langle r\rangle\) (тест 33); рис. 12 --- дираковский
конус \(E_{\min}(1/L)\); рис. 13 --- дираковский провал DOS; рис. 14 ---
Байерс--Янг (целый против дробного заряда); рис. 15 --- Берри-обрез
\(R_2(0;T)\); рис. 16 --- Хатано--Нельсон, эллипс неэрмитового спектра;
рис. 17 --- спинорная фаза \(\gamma^{\ast}\); рис. 18 --- фазовая
текстура вихревой пары; рис. 19 --- бабочка Хофштадтера с отмеченной
точкой \(\alpha=1/2\).

\section{Приложение D. Верификация всех 64 спинорных структур квартики
Клейна (v22.1,
spinor64)}\label{ux43fux440ux438ux43bux43eux436ux435ux43dux438ux435-d.-ux432ux435ux440ux438ux444ux438ux43aux430ux446ux438ux44f-ux432ux441ux435ux445-64-ux441ux43fux438ux43dux43eux440ux43dux44bux445-ux441ux442ux440ux443ux43aux442ux443ux440-ux43aux432ux430ux440ux442ux438ux43aux438-ux43aux43bux435ux439ux43dux430-v22.1-spinor64}

Эталонная реализация \texttt{verification/spinor64} (Python/NumPy,
воспроизведение:
\texttt{python3\ verification/spinor64/run\_spinor64.py}) выполняет два
независимых эксперимента по всем 64 спинорным структурам
\(\varepsilon\in\mathbb{F}_2^6\) квартики Клейна и исправляет
утверждение монографии v21 о «уникальности» структуры idx=38.

\textbf{E1 --- точная симметрия (граф Клейна \{3,7\}).} Квартика Клейна
дискретизована своим регулярным отображением \{3,7\} (56 вершин, 84
ребра, 24 семиугольника, Aut = PSL(2,7) порядка 168). 64 спин-структуры
реализованы как сигнатуры рёбер с нечётной граневой чётностью (модель
Кастелейна--Чимасони--Решетихина, каноническая калибровка: все рёбра
остовного дерева положительны). Численно найдено разбиение
PSL(2,7)-действия на орбиты \textbf{28 / 21 / 7 / 7 / 1}; орбита из 28
--- нечётные структуры (Arf=1) --- единственная, что подтверждает
транзитивность PSL(2,7) на 28 битангентах (теорема Римана--Клейна).
Знаковые операторы Дирака внутри орбиты перестановочно-сопряжены:
максимальное попарное расстояние между спектрами по всем 64 структурам
\textbf{8.9·10⁻¹⁵}; калибровочная инвариантность 7.1·10⁻¹⁵; число
нулевых мод дискретного оператора: 2 (нечётная орбита) / 3 (чётные
орбиты) / 7 (тривиальный класс). Следствие: никакая спинорная структура
не может обладать уникальной статистикой --- «уникальность idx=38» в v21
была артефактом дискретизации, разрушающим симметрию PSL(2,7) (см. разд.
3.2.3 оитога v21.1). По собственной формуле v21
\(\mathrm{Arf}(\varepsilon)=
\varepsilon_1\varepsilon_2+\varepsilon_3\varepsilon_4+\varepsilon_5\varepsilon_6\)
вектор \(\varepsilon(38)=(0,1,1,0,0,1)\) даёт Arf = 0 (а не 1, как
утверждалось).

\textbf{E2 --- статистика (AB-облако).} Хофштадтеровский тор \(L=44\),
\(\alpha=1/2\), \(N_v=54\) вихрей \(q=+1\) (плотностное масштабирование
якоря \(N_v=25\) при \(30\times30\)), \(W=0\), калибровка
\texttt{:monumental} (atan-гладкая, вертикальные связи, фактор 0.5),
спин-структура --- граничные твисты
\(\varphi_x=\pi(\varepsilon_1+\varepsilon_3+\varepsilon_5)\),
\(\varphi_y=\pi(\varepsilon_2+\varepsilon_4+\varepsilon_6)\); усреднение
по 5 конфигурациям вихрей, объёмное окно 0.6. Референс --- Monte-Carlo
GUE ансамбль 100 матриц \(1936\times1936\) (методология теста 16):
медиана \(\langle r\rangle = 0.6013\), 95\% ДИ \([0.5847; 0.6140]\).
Результат: \textbf{все 64 структуры внутри ДИ}:
\(\langle r\rangle = 0.5984\pm0.0035\) (разброс \(0.5935\dots0.6027\),
min MC p = 0.36). GUE-статистика одинакова для всех структур ---
источник GUE, как и заключено в разделе 4.1, --- динамика AB-облака, а
не выбор спинорной структуры.

Порты теста 38 (замороженные данные + собственный алгоритм Якоби, без
LAPACK) поставляются на 10 языках:
\texttt{verification/\{cpp,java,rust,go,fortran,haskell,r,matlab,julia,javascript\}/spinor38/}.

\end{document}
